\documentclass[10pt,a4paper,twocolumn]{article}
\usepackage{subfiles}

% ==============================================================================
% banq-preamble.tex
% Shared LaTeX preamble for all XPower Banq papers and the unified bundle.
% Loaded via % ==============================================================================
% banq-preamble.tex
% Shared LaTeX preamble for all XPower Banq papers and the unified bundle.
% Loaded via % ==============================================================================
% banq-preamble.tex
% Shared LaTeX preamble for all XPower Banq papers and the unified bundle.
% Loaded via \input{../shared/banq-preamble} from each child paper and from
% the master bundle root. Identical to the original per-paper preambles.
% ==============================================================================

\usepackage[utf8]{inputenc}
\usepackage[T1]{fontenc}
\usepackage{newunicodechar}
\usepackage{libertinus}
\usepackage[margin=2cm]{geometry}
\usepackage{amsmath,amssymb,amsthm}
\usepackage{graphicx}
\usepackage{xcolor}
\usepackage[hyphens]{url}
\usepackage{hyperref}
\usepackage{booktabs}
\usepackage{tabularx}
\usepackage{multirow}
\usepackage{enumitem}
\usepackage{tikz}
\usepackage{pgfplots}
\usepackage{listings}
\usepackage{caption}
\usepackage{subcaption}
\usepackage{fancyhdr}
\usepackage{float}
\usepackage{placeins}  % For \FloatBarrier
\usepackage{tocloft}   % Adjustable TOC number-box widths
\usepackage{multicol}  % Multi-column layout (used for the bundle's master TOC)
\usepackage{twemojis}  % Twemoji glyphs (pdflatex-compatible color emoji)

% Float parameters to prevent overlapping
\setcounter{topnumber}{2}
\setcounter{bottomnumber}{2}
\setcounter{totalnumber}{4}
\renewcommand{\topfraction}{0.85}
\renewcommand{\bottomfraction}{0.85}
\renewcommand{\textfraction}{0.15}
\renewcommand{\floatpagefraction}{0.7}

% Prevent line breaks inside inline math (fixes broken superscripts/sqrt)
\relpenalty=10000
\binoppenalty=10000

\pgfplotsset{compat=1.18}
\usetikzlibrary{shapes,arrows,positioning,calc,patterns}

% ==============================================================================
% Colors and Styling
% ==============================================================================
\definecolor{codeblue}{RGB}{0,102,204}
\definecolor{codepurple}{RGB}{128,0,128}
\definecolor{codegreen}{RGB}{0,128,0}
\definecolor{codegray}{RGB}{128,128,128}
\definecolor{primaryblue}{RGB}{41,98,255}
\definecolor{warningred}{RGB}{220,53,69}
\definecolor{successgreen}{RGB}{40,167,69}

\hypersetup{
    colorlinks=true,
    linkcolor=primaryblue,
    citecolor=primaryblue,
    urlcolor=primaryblue
}

\lstset{
    basicstyle=\ttfamily\scriptsize,
    keywordstyle=\color{codeblue}\bfseries,
    commentstyle=\color{codegreen},
    stringstyle=\color{codepurple},
    numbers=left,
    numberstyle=\tiny\color{codegray},
    numbersep=4pt,
    xleftmargin=2em,
    framexleftmargin=2em,
    breaklines=true,
    frame=single,
    captionpos=b
}
\lstdefinelanguage{Solidity}{
    keywords={function, returns, return, if, else, for, while, do, break, continue, throw, require, revert, assert, public, private, internal, external, view, pure, payable, memory, storage, calldata, constant, immutable, virtual, override, modifier, event, emit, struct, enum, mapping, contract, interface, library, abstract, is, new, delete, true, false, uint256, uint128, uint64, uint32, uint16, uint8, uint240, uint224, int256, int128, int64, int32, int8, address, bool, bytes, bytes32, string, unchecked},
    sensitive=true,
    morecomment=[l]{//},
    morecomment=[s]{/*}{*/},
    morestring=[b]",
    morestring=[b]'
}

% ==============================================================================
% Theorem Environments
% ==============================================================================
\theoremstyle{definition}
\newtheorem{definition}{Definition}[section]
\newtheorem{property}{Property}[section]
\theoremstyle{plain}
\newtheorem{theorem}{Theorem}[section]
\newtheorem{lemma}[theorem]{Lemma}
\newtheorem{corollary}[theorem]{Corollary}
\theoremstyle{remark}
\newtheorem{remark}{Remark}[section]
 from each child paper and from
% the master bundle root. Identical to the original per-paper preambles.
% ==============================================================================

\usepackage[utf8]{inputenc}
\usepackage[T1]{fontenc}
\usepackage{newunicodechar}
\usepackage{libertinus}
\usepackage[margin=2cm]{geometry}
\usepackage{amsmath,amssymb,amsthm}
\usepackage{graphicx}
\usepackage{xcolor}
\usepackage[hyphens]{url}
\usepackage{hyperref}
\usepackage{booktabs}
\usepackage{tabularx}
\usepackage{multirow}
\usepackage{enumitem}
\usepackage{tikz}
\usepackage{pgfplots}
\usepackage{listings}
\usepackage{caption}
\usepackage{subcaption}
\usepackage{fancyhdr}
\usepackage{float}
\usepackage{placeins}  % For \FloatBarrier
\usepackage{tocloft}   % Adjustable TOC number-box widths
\usepackage{multicol}  % Multi-column layout (used for the bundle's master TOC)
\usepackage{twemojis}  % Twemoji glyphs (pdflatex-compatible color emoji)

% Float parameters to prevent overlapping
\setcounter{topnumber}{2}
\setcounter{bottomnumber}{2}
\setcounter{totalnumber}{4}
\renewcommand{\topfraction}{0.85}
\renewcommand{\bottomfraction}{0.85}
\renewcommand{\textfraction}{0.15}
\renewcommand{\floatpagefraction}{0.7}

% Prevent line breaks inside inline math (fixes broken superscripts/sqrt)
\relpenalty=10000
\binoppenalty=10000

\pgfplotsset{compat=1.18}
\usetikzlibrary{shapes,arrows,positioning,calc,patterns}

% ==============================================================================
% Colors and Styling
% ==============================================================================
\definecolor{codeblue}{RGB}{0,102,204}
\definecolor{codepurple}{RGB}{128,0,128}
\definecolor{codegreen}{RGB}{0,128,0}
\definecolor{codegray}{RGB}{128,128,128}
\definecolor{primaryblue}{RGB}{41,98,255}
\definecolor{warningred}{RGB}{220,53,69}
\definecolor{successgreen}{RGB}{40,167,69}

\hypersetup{
    colorlinks=true,
    linkcolor=primaryblue,
    citecolor=primaryblue,
    urlcolor=primaryblue
}

\lstset{
    basicstyle=\ttfamily\scriptsize,
    keywordstyle=\color{codeblue}\bfseries,
    commentstyle=\color{codegreen},
    stringstyle=\color{codepurple},
    numbers=left,
    numberstyle=\tiny\color{codegray},
    numbersep=4pt,
    xleftmargin=2em,
    framexleftmargin=2em,
    breaklines=true,
    frame=single,
    captionpos=b
}
\lstdefinelanguage{Solidity}{
    keywords={function, returns, return, if, else, for, while, do, break, continue, throw, require, revert, assert, public, private, internal, external, view, pure, payable, memory, storage, calldata, constant, immutable, virtual, override, modifier, event, emit, struct, enum, mapping, contract, interface, library, abstract, is, new, delete, true, false, uint256, uint128, uint64, uint32, uint16, uint8, uint240, uint224, int256, int128, int64, int32, int8, address, bool, bytes, bytes32, string, unchecked},
    sensitive=true,
    morecomment=[l]{//},
    morecomment=[s]{/*}{*/},
    morestring=[b]",
    morestring=[b]'
}

% ==============================================================================
% Theorem Environments
% ==============================================================================
\theoremstyle{definition}
\newtheorem{definition}{Definition}[section]
\newtheorem{property}{Property}[section]
\theoremstyle{plain}
\newtheorem{theorem}{Theorem}[section]
\newtheorem{lemma}[theorem]{Lemma}
\newtheorem{corollary}[theorem]{Corollary}
\theoremstyle{remark}
\newtheorem{remark}{Remark}[section]
 from each child paper and from
% the master bundle root. Identical to the original per-paper preambles.
% ==============================================================================

\usepackage[utf8]{inputenc}
\usepackage[T1]{fontenc}
\usepackage{newunicodechar}
\usepackage{libertinus}
\usepackage[margin=2cm]{geometry}
\usepackage{amsmath,amssymb,amsthm}
\usepackage{graphicx}
\usepackage{xcolor}
\usepackage[hyphens]{url}
\usepackage{hyperref}
\usepackage{booktabs}
\usepackage{tabularx}
\usepackage{multirow}
\usepackage{enumitem}
\usepackage{tikz}
\usepackage{pgfplots}
\usepackage{listings}
\usepackage{caption}
\usepackage{subcaption}
\usepackage{fancyhdr}
\usepackage{float}
\usepackage{placeins}  % For \FloatBarrier
\usepackage{tocloft}   % Adjustable TOC number-box widths
\usepackage{multicol}  % Multi-column layout (used for the bundle's master TOC)
\usepackage{twemojis}  % Twemoji glyphs (pdflatex-compatible color emoji)

% Float parameters to prevent overlapping
\setcounter{topnumber}{2}
\setcounter{bottomnumber}{2}
\setcounter{totalnumber}{4}
\renewcommand{\topfraction}{0.85}
\renewcommand{\bottomfraction}{0.85}
\renewcommand{\textfraction}{0.15}
\renewcommand{\floatpagefraction}{0.7}

% Prevent line breaks inside inline math (fixes broken superscripts/sqrt)
\relpenalty=10000
\binoppenalty=10000

\pgfplotsset{compat=1.18}
\usetikzlibrary{shapes,arrows,positioning,calc,patterns}

% ==============================================================================
% Colors and Styling
% ==============================================================================
\definecolor{codeblue}{RGB}{0,102,204}
\definecolor{codepurple}{RGB}{128,0,128}
\definecolor{codegreen}{RGB}{0,128,0}
\definecolor{codegray}{RGB}{128,128,128}
\definecolor{primaryblue}{RGB}{41,98,255}
\definecolor{warningred}{RGB}{220,53,69}
\definecolor{successgreen}{RGB}{40,167,69}

\hypersetup{
    colorlinks=true,
    linkcolor=primaryblue,
    citecolor=primaryblue,
    urlcolor=primaryblue
}

\lstset{
    basicstyle=\ttfamily\scriptsize,
    keywordstyle=\color{codeblue}\bfseries,
    commentstyle=\color{codegreen},
    stringstyle=\color{codepurple},
    numbers=left,
    numberstyle=\tiny\color{codegray},
    numbersep=4pt,
    xleftmargin=2em,
    framexleftmargin=2em,
    breaklines=true,
    frame=single,
    captionpos=b
}
\lstdefinelanguage{Solidity}{
    keywords={function, returns, return, if, else, for, while, do, break, continue, throw, require, revert, assert, public, private, internal, external, view, pure, payable, memory, storage, calldata, constant, immutable, virtual, override, modifier, event, emit, struct, enum, mapping, contract, interface, library, abstract, is, new, delete, true, false, uint256, uint128, uint64, uint32, uint16, uint8, uint240, uint224, int256, int128, int64, int32, int8, address, bool, bytes, bytes32, string, unchecked},
    sensitive=true,
    morecomment=[l]{//},
    morecomment=[s]{/*}{*/},
    morestring=[b]",
    morestring=[b]'
}

% ==============================================================================
% Theorem Environments
% ==============================================================================
\theoremstyle{definition}
\newtheorem{definition}{Definition}[section]
\newtheorem{property}{Property}[section]
\theoremstyle{plain}
\newtheorem{theorem}{Theorem}[section]
\newtheorem{lemma}[theorem]{Lemma}
\newtheorem{corollary}[theorem]{Corollary}
\theoremstyle{remark}
\newtheorem{remark}{Remark}[section]

% ==============================================================================
% banq-meta.tex
% Shared author/version metadata for all XPower Banq papers.
% ==============================================================================

\newcommand{\banqAuthor}{\textsc{Karun The Ritch}}
\newcommand{\banqAuthorEmail}{ktr@xpowerbanq.com}
\newcommand{\banqAuthorURL}{https://www.xpowerbanq.com}
\newcommand{\banqAuthorContact}{%
    \href{mailto:\banqAuthorEmail}{\texttt{\banqAuthorEmail}}\quad\url{\banqAuthorURL}%
}
\newcommand{\banqArxivId}{preprint:2605~~[q-fin.CP]}
\newcommand{\banqArxivDate}{1 May 2026}

% ==============================================================================
% banq-bundle.tex
% Defines the \ifbundle conditional that distinguishes a standalone child build
% from the unified bundle build, and the \xref macro that produces a real
% internal hyperlink in bundle mode and a citation-with-text in standalone mode.
%
% Standalone child papers leave \bundlefalse (the default).
% The master bundle root sets \bundletrue before \begin{document}.
% ==============================================================================

\newif\ifbundle
\bundlefalse

% \xref{<bibkey>}{<label>}{<text>}
% - In bundle mode: produces a clickable in-document hyperlink to the label.
% - In standalone mode: prints the text followed by a citation to the bibkey.
% Note: defined after \bundletrue/\bundlefalse so it picks up the right branch.
% Children must call this from inside \AtBeginDocument to ensure \ifbundle is
% finalized; alternatively the macro can test \ifbundle at call site.
\newcommand{\xref}[3]{%
  \ifbundle
    \hyperref[#2]{#3}%
  \else
    #3~\cite{#1}%
  \fi
}

% \xrefshort{<bibkey>}{<label>}{<text>}
% Same as \xref but suppresses the citation entirely in standalone mode (used
% when the surrounding prose already cites the source).
\newcommand{\xrefshort}[3]{%
  \ifbundle
    \hyperref[#2]{#3}%
  \else
    #3%
  \fi
}

% \banqEnableBundleNumbering
% Resets the section counter at every \part and prefixes \thesection with the
% part's Roman numeral, so headings read "I.1 Introduction", "II.1 Introduction",
% etc. Called explicitly from the master root after \bundletrue (see
% 000.banq-all/banq.tex). Not called by standalone children.
%
% Parts that contain MULTIPLE sub-papers must additionally call \banqSubpart
% (defined below) before each child \subfile, which appends a per-paper letter
% suffix (a, b, ...) to the part numeral and resets the section counter so each
% sub-paper restarts from 1. Without this, sub-papers within the same part
% share a single section sequence, which causes the TOC to conflate identically
% titled sections (e.g. two "II.1 Introduction" entries for banq-lck/banq-log).
%
% The outer \makeatletter/\makeatother brackets here are critical: the macro
% body contains \@addtoreset (an @ control sequence) which must be parseable
% as a single token at \newcommand-definition time, not at call time.
\newcommand{\banqSubpartLetter}{}
% Boolean tracking whether we are currently inside a sub-paper of a shared part
% (i.e. \banqSubpart has been called since the last \part). Used by \banqPaper
% to decide whether to demote section bookmark levels. \ifx-comparing
% \banqSubpartLetter against \@empty doesn't work because \newcommand produces
% \long macros while \@empty is short — \ifx always returns false.
\newif\ifbanqSharedPart
\makeatletter
\newcommand{\banqEnableBundleNumbering}{%
  \@addtoreset{section}{part}%
  \renewcommand{\thesection}{\Roman{part}\banqSubpartLetter.\,\arabic{section}}%
  % Make hyperref anchor names unique across sub-papers that share a \part.
  % Without this, banq-lck and banq-log both emit section.2.1, section.2.2, …
  % which corrupts the PDF outline tree (external viewers truncate at Part II).
  \renewcommand{\theHsection}{\Roman{part}\banqSubpartLetter.\arabic{section}}%
  \renewcommand{\theHsubsection}{\theHsection.\arabic{subsection}}%
  \renewcommand{\theHsubsubsection}{\theHsubsection.\arabic{subsubsection}}%
  \renewcommand{\theHparagraph}{\theHsubsubsection.\arabic{paragraph}}%
  % Hook \part so the subpart-letter and shared-part flag both reset when a new
  % part begins; otherwise state set in part N would leak into part N+1.
  \let\banqOriginalPart\part
  \renewcommand{\part}{%
    \renewcommand{\banqSubpartLetter}{}%
    \banqSharedPartfalse
    \banqOriginalPart
  }%
  % The "I.\,10" prefix is wider than the default \numberline box. Widen the
  % TOC number-box widths so two-digit section numbers don't overlap titles.
  \setlength{\cftsecnumwidth}{3.5em}%
  \setlength{\cftsubsecnumwidth}{4.2em}%
  \setlength{\cftsubsubsecnumwidth}{5.2em}%
}
\makeatother

% \banqSubpart{<letter>}
% Marks the start of a sub-paper inside a part that contains multiple papers.
% Sets the per-paper letter suffix appended to the part numeral in \thesection
% (so headings read "IIa.1 Introduction", "IIb.1 Introduction", etc.) and
% resets the section counter so the sub-paper's own numbering restarts at 1.
% Pass an empty argument (\banqSubpart{}) to clear the suffix mid-part if
% needed; the suffix also auto-clears at every \part boundary (see hook above).
\newcommand{\banqSubpart}[1]{%
  \renewcommand{\banqSubpartLetter}{#1}%
  \setcounter{section}{0}%
  \banqSharedParttrue
}

% \banq@RestoreLevels / \banq@DemoteLevels
% Adjust hyperref's per-sectioning bookmark levels. When two sub-papers share a
% \part (banq-lck "IIa" and banq-log "IIb" inside Part~II), each paper's
% sections should nest under its own divider bookmark in the PDF outline,
% rather than appearing as siblings of it under the part. Demoting
% \toclevel@section from 1 to 2 (and the deeper levels accordingly) achieves
% that nesting; \banqPaper calls Restore before emitting its own divider so the
% divider always lands at level 1, then calls Demote afterwards (only when in a
% shared part) so subsequent sections sit at level 2 = child of the divider.
% Affects only the PDF outline depths; the printed TOC layout is unchanged.
\makeatletter
\newcommand{\banq@RestoreLevels}{%
  \renewcommand{\toclevel@section}{1}%
  \renewcommand{\toclevel@subsection}{2}%
  \renewcommand{\toclevel@subsubsection}{3}%
  \renewcommand{\toclevel@paragraph}{4}%
}
\newcommand{\banq@DemoteLevels}{%
  \renewcommand{\toclevel@section}{2}%
  \renewcommand{\toclevel@subsection}{3}%
  \renewcommand{\toclevel@subsubsection}{4}%
  \renewcommand{\toclevel@paragraph}{5}%
}
\makeatother

% Storage for the current part's twemoji name. \banqPartWithEmoji sets it
% (overwritten at each part boundary). \banqPaper reads it without clearing,
% so multiple sub-papers within the same part (e.g. banq-lck and banq-log
% in Part~II) share the same emoji on their respective cover pages.
\makeatletter
\def\banq@partemoji{}
\newif\ifbanq@hasemoji
\banq@hasemojifalse
% \banqSetSubpartEmoji{<twemoji-name>}
% Override the part's default emoji for the next sub-paper cover. Used inside
% parts that contain multiple papers when each sub-paper deserves its own
% cover icon (e.g. banq-lck takes "lock" while banq-log keeps the part-level
% "gear"). Encapsulates the @-internal \gdef so the master document body need
% not enter \makeatletter scope.
\newcommand{\banqSetSubpartEmoji}[1]{\gdef\banq@partemoji{#1}}
\makeatother

% \banqForceRecto
% Page-break helper that flushes outstanding output and ensures the next
% content-bearing page lands on an odd (recto) page in duplex bindings.
% The article class is loaded without `twoside` so \cleardoublepage degenerates
% to \clearpage and will not pad; we check the page counter explicitly and
% insert a blank filler when the upcoming page is even. \value{page} reflects
% the next page number after \clearpage, before any content has been emitted
% on it. Call sites: front-matter (TOC start, body start) and \banqPaper.
\newcommand{\banqForceRecto}{%
  \clearpage
  \ifodd\value{page}\else\hbox{}\thispagestyle{empty}\newpage\fi
}

% \banqEnsureEvenPageCount
% Final-pass parity helper: guarantees the document ends with an even total
% page count, so duplex-bound copies have a back cover on the verso side and
% any subsequent insert (e.g. a colophon, a bound-in errata sheet) starts on
% a recto. After \clearpage, \value{page} holds the *next* page number, so
% the current page count equals \value{page}-1; the count is even iff
% \value{page} is odd, which is exactly the invariant \banqForceRecto already
% enforces. The two operations are semantically distinct (recto-start vs.
% even-total) but compute identical pad logic, so we alias them. Call once,
% immediately before \end{document}, in the bundle root.
\let\banqEnsureEvenPageCount\banqForceRecto

% \banqPartWithEmoji{<title>}{<twemoji-name>}
% Bumps the part counter, stashes the title for the part-level TOC entry,
% and stashes the emoji for the cover page(s) that follow (see \banqPaper).
% No physical divider page is emitted — the cover that immediately follows
% is the visible boundary between parts. The bundle's section-numbering
% scheme (I.1, IIa.1, IIb.1, III.1, ...) continues to work because
% \refstepcounter drives the same counter \part would have, and
% \@addtoreset (set up by \banqEnableBundleNumbering) still resets section
% at the part boundary.
%
% The TOC entry itself is written by \banqPaper at the start of the cover
% page rather than here, so the deferred page-number capture lands on the
% cover (the visible part boundary) rather than on whatever page the build
% happens to be on between subfiles.
\makeatletter
\def\banq@parttitle{}
\newif\ifbanq@parttocpending
\banq@parttocpendingfalse
\newcommand{\banqPartWithEmoji}[2]{%
  \refstepcounter{part}%
  \renewcommand{\banqSubpartLetter}{}%
  \banqSharedPartfalse
  \gdef\banq@partemoji{#2}%
  \global\banq@hasemojitrue
  \gdef\banq@parttitle{#1}%
  \global\banq@parttocpendingtrue
}
\makeatother

% \banqAbstract{<abstract text>}{<keywords>}
% Single source of truth for each child's abstract + keywords list, with
% mode-aware rendering:
%   - Standalone: emits a normal \begin{abstract}...\end{abstract} block
%     followed by a "Keywords:" line (the previous \ifbundle\else ... \fi
%     boilerplate each child used to inline).
%   - Bundle: stashes the content into module-level token lists and sets
%     \ifbanq@hasabstract; \banqPaper then renders it on the verso page that
%     follows the cover (the page that was previously left blank).
% Defined as \long so abstracts containing paragraph breaks are accepted,
% though in practice all six abstracts are single-paragraph.
\makeatletter
\def\banq@abstracttext{}
\def\banq@keywordstext{}
\newif\ifbanq@hasabstract
\banq@hasabstractfalse
\long\def\banqAbstract#1#2{%
  \ifbundle
    \long\gdef\banq@abstracttext{#1}%
    \long\gdef\banq@keywordstext{#2}%
    \global\banq@hasabstracttrue
  \else
    \begin{abstract}#1\end{abstract}\par\medskip\noindent
    \textbf{Keywords:} #2\par
  \fi
}
\makeatother

% \banqPaper{<title>}{<subtitle>}
% Paper-level divider page printed at the top of each child's body when in
% bundle mode. Sits below the \part page in the visual hierarchy and provides
% a clear transition between consecutive papers within the same part (e.g.
% banq-lck -> banq-log inside Part~II). A no-op in standalone builds.
%
% Wrapped in \makeatletter so the macro body can reference \banq@RestoreLevels,
% \banq@DemoteLevels, and \@empty as single tokens rather than \banq + @… .
\makeatletter
\newcommand{\banqPaper}[2]{%
  \ifbundle
    % Force the cover page to land on an odd (recto) page; see \banqForceRecto.
    \banqForceRecto
    \onecolumn
    \thispagestyle{empty}
    % If a \banqPartWithEmoji call is pending, write its TOC + PDF-outline
    % entry now so the captured page number is the cover page (the visible
    % part boundary), not whatever blank pad the build happened to land on.
    \ifbanq@parttocpending
      \phantomsection
      \addcontentsline{toc}{part}{\thepart\hspace{1em}\banq@parttitle}%
      \global\banq@parttocpendingfalse
    \fi
    \vspace*{6em}
    \begin{center}
      {\Large\color{codegray} \textsc{XPower Banq: Part \Roman{part}\banqSubpartLetter}}\\[0.6em]
      \rule{0.4\textwidth}{0.4pt}\\[1.2em]
      {\fontsize{22}{26}\selectfont\bfseries #1}\\[0.8em]
      {\large\color{codegray} #2}
    \end{center}
    \vspace*{\fill}
    \ifbanq@hasemoji
      \begin{center}
        \twemoji[height=8cm]{\banq@partemoji}
      \end{center}
    \fi
    \vspace*{\fill}
    \clearpage
    % Verso page after the cover. If the child supplied an abstract via
    % \banqAbstract, render it here (still in onecolumn, unstyled body width
    % so the prose breathes). Otherwise leave the page blank — used for the
    % "References & Glossary" cover, which has no per-paper abstract since
    % the bundle's front-matter abstract already summarises the whole volume.
    \thispagestyle{empty}
    \ifbanq@hasabstract
      \vspace*{\fill}
      \begin{center}{\large\bfseries Abstract}\end{center}
      \medskip
      \begingroup
        \small
        \noindent\banq@abstracttext\par
        \medskip
        \noindent\textbf{Keywords:} \banq@keywordstext\par
      \endgroup
      \vspace*{\fill}
      \global\banq@hasabstractfalse
    \else
      \hbox{}%
    \fi
    \newpage
    \twocolumn
    % Register the divider in the TOC and PDF outline as a child of the part.
    % Restore section bookmark levels first so this divider always lands at
    % level 1 (regardless of any prior demotion left over from a sibling
    % sub-paper), then conditionally demote so subsequent sections nest under
    % this divider when we're inside a shared part.
    \banq@RestoreLevels
    \phantomsection
    \addcontentsline{toc}{section}{\protect\numberline{\Roman{part}\banqSubpartLetter.}#1}%
    \ifbanqSharedPart
      \banq@DemoteLevels
    \fi
  \fi
}
\makeatother


\title{\textbf{\Large XPower Banq: Simulations \& Risk Analysis}\\[0.3em]
\normalsize Capacity Accumulation, Cascade, TWAP, and Bad-Debt Simulations\\
for the XPower Banq Lending Protocol}

\author{
    \banqAuthor\thanks{\banqAuthorContact}
}

\date{}

\ifbundle\else
\renewcommand{\thepage}{SIM/\arabic{page}}
\fi

\begin{document}

\ifbundle\else
\maketitle
% ==============================================================================
% arxiv-tag.tex
% The arXiv-style vertical preprint tag printed on the first page of each
% standalone paper and on the bundle title page. The date is taken from
% \banqArxivDate (defined in banq-meta.tex) and may be overridden locally
% before % ==============================================================================
% arxiv-tag.tex
% The arXiv-style vertical preprint tag printed on the first page of each
% standalone paper and on the bundle title page. The date is taken from
% \banqArxivDate (defined in banq-meta.tex) and may be overridden locally
% before % ==============================================================================
% arxiv-tag.tex
% The arXiv-style vertical preprint tag printed on the first page of each
% standalone paper and on the bundle title page. The date is taken from
% \banqArxivDate (defined in banq-meta.tex) and may be overridden locally
% before \input{../shared/arxiv-tag} via \renewcommand.
% ==============================================================================

\begin{tikzpicture}[remember picture, overlay]
  \node[rotate=90, anchor=center, font=\Huge, color=codegray]
    at ([xshift=1cm]current page.west)
    {\banqArxivId~~\banqArxivDate};
\end{tikzpicture}
 via \renewcommand.
% ==============================================================================

\begin{tikzpicture}[remember picture, overlay]
  \node[rotate=90, anchor=center, font=\Huge, color=codegray]
    at ([xshift=1cm]current page.west)
    {\banqArxivId~~\banqArxivDate};
\end{tikzpicture}
 via \renewcommand.
% ==============================================================================

\begin{tikzpicture}[remember picture, overlay]
  \node[rotate=90, anchor=center, font=\Huge, color=codegray]
    at ([xshift=1cm]current page.west)
    {\banqArxivId~~\banqArxivDate};
\end{tikzpicture}

\fi

\banqAbstract{%
We present the simulation suite that empirically validates the XPower Banq lending protocol~\cite{wp}. Four studies are reported: (1)~capacity accumulation under the beta-distributed cap function, confirming $O(\sqrt{n})$ Sybil-resistance scaling; (2)~liquidation-cascade simulations across five market-impact regimes, demonstrating up to 80\% cascade reduction at full lock adoption; (3)~log-space TWAP oracle response to price shocks and liquidity changes, including 60 on-chain Foundry scenario tests showing two-tick manipulation immunity; and (4)~bad-debt risk quantification via Merton jump-diffusion Monte Carlo, with closed-form analytical bound and a five-configuration safe operating region. The mathematical foundations underlying these simulations are developed in the companion theory paper~\cite{mtp}.%
}{DeFi simulation, Monte Carlo, jump diffusion, TWAP, cascade dynamics, bad debt, Sybil resistance}

\banqPaper{Simulations \& Risk Analysis}{Capacity Accumulation, Cascade, TWAP, and Bad-Debt Simulations for the XPower Banq Lending Protocol}

\section{Capacity Accumulation Simulation}
\label{apx:simulation}

The following Python script simulates capacity accumulation under the beta-distributed cap function (the beta cap formula (\S4.4 of~\cite{wp})).

\begin{lstlisting}[language=Python, caption={Capacity accumulation simulation}, label=lst:simulate]
#!/usr/bin/env python3
from numpy import sqrt

def my_cap(balance, max_cap, supply, n_accounts) -> float:
    """
    Computes the max. individual capacity gain per round
    for given balance, supply, max. capacity and accounts.
    """
    # compute lambda: balance / supply
    lmb = balance / supply
    # compute my capacity gain multiplier
    mul = 12*lmb*(1-lmb)**2
    # compute my capacity gain divisor
    div = sqrt(n_accounts + 2)
    # compute my capacity gain
    return max_cap * mul / div

if __name__ == "__main__":
    # epsilon for stopping criterion
    eps = 1e-3
    # max. capacity (normalized)
    max_cap = 1.00
    # supply (normalized)
    supply = 1.00

    print(f"beg_balance;n_accounts;iterations")
    for beg_balance in map(lambda e: 10**(-e), range(1,9)):
        for n_accounts in map(lambda e: 10**(2*e), range(1,5)):

            end_balance = beg_balance
            iterations = 0

            # simulate cap gain until close to supply
            while (supply - end_balance) / supply > eps:
                end_balance += my_cap(
                    end_balance, max_cap, supply, n_accounts
                )
                iterations += 1

            print(f"{beg_balance:.0e};{n_accounts:.0e};{iterations:.3e}")
\end{lstlisting}


\subsection{Cap Function Components}

The \texttt{my\_cap} function implements the beta cap formula (\S4.4 of~\cite{wp}). The balance ratio $\lambda = B/S$ normalizes across token amounts and decimals. The Beta(2,3) multiplier $12\lambda(1{-}\lambda)^2$ vanishes at both boundaries ($\lambda = 0$ and $\lambda = 1$), peaks at $\lambda = 1/3$, and decays quadratically via $(1{-}\lambda)^2$ as holdings approach monopoly. The holder divisor $\sqrt{n{+}2}$ provides sublinear Sybil resistance: creating more accounts increases $n$, reducing per-account cap gains, with the offset of~2 preventing degeneracy at low holder counts. $C_{\max}$ is the governance-configured upper bound on relative capacity. In the protocol, \texttt{max\_cap} is the residual capacity \texttt{\_relLimit(abs\_limit, total)}; the simulation's normalised \texttt{max\_cap = 1.00, supply = 1.00} corresponds to the limiting case where \texttt{abs\_limit} $\gg$ \texttt{total} so the headroom is approximately \texttt{abs\_limit}.

\subsection{Simulation Algorithm}

The simulation sweeps over starting balances and holder counts. For each combination, it iteratively adds the cap gain (via \texttt{my\_cap}) to the user's balance until the remaining capacity $(1 - \lambda)$ falls below $\varepsilon = 0.1\%$ of total supply, recording the number of iterations required. The stopping threshold avoids infinite loops caused by the asymptotic decay of the cap function as $\lambda \to 1$. The holder count $n$ is held constant throughout each run; in practice, $n$ would fluctuate as users enter and exit.

\subsection{Convergence Behavior}

Near $\lambda = 0$, the cap gain grows linearly (${\approx}\, 12\lambda / \sqrt{n{+}2}$), producing a fast start. Near $\lambda = 1$, it shrinks quadratically (${\approx}\, 12(1{-}\lambda)^2 / \sqrt{n{+}2}$), producing a slow approach to full capacity. Because most iterations occur in this slow-convergence region, starting balance has minimal impact on total iterations. Since cap gain scales as $1/\sqrt{n}$, iterations scale as $\sqrt{n}$---explaining the $10\times$ increase when moving from $10^2$ to $10^4$ accounts.

\subsection{Representative Output}

\begin{table}[H]
\centering
\caption{Simulation results: iterations to full capacity}
\label{tab:simulation-output}
\small
\begin{tabularx}{0.95\columnwidth}{@{}>{\centering\arraybackslash}X >{\centering\arraybackslash}X >{\raggedleft\arraybackslash}X@{}}
\toprule
Starting $\lambda$ & Accounts ($n$) & Iterations \\
\midrule
$10^{-2}$ & $10^2$ & 846 \\
$10^{-2}$ & $10^4$ & 8,417 \\
$10^{-2}$ & $10^6$ & 84,200 \\
$10^{-2}$ & $10^8$ & 842,100 \\
\midrule
$10^{-4}$ & $10^2$ & 852 \\
$10^{-4}$ & $10^4$ & 8,458 \\
$10^{-4}$ & $10^6$ & 84,590 \\
$10^{-4}$ & $10^8$ & 845,900 \\
\bottomrule
\end{tabularx}
\end{table}

The results confirm the $O(\sqrt{n})$ scaling: $10^4$ accounts require approximately $10\times$ more iterations than $10^2$ accounts. Starting balance has minimal impact ($<$2\% across 7 orders of magnitude), confirming that convergence near $\lambda = 1$ dominates. For a protocol with $10^6$ large holders, any user requires approximately 84,000 iterations to reach full capacity, providing predictable growth independent of initial position size.

\subsection{Multi-Account Share Dynamics}

The following simulation examines how capital shares evolve when multiple accounts compete for capacity under the beta cap function and iteration constraints:

\begin{lstlisting}[language=Python, caption={Multi-account share dynamics simulation}, label=lst:share-sim]
#!/usr/bin/env python3
from numpy import sqrt

def beta_cap(balance, max_cap, supply, n_holders):
    """Beta-distributed cap function (Equation 4)."""
    lam = balance / supply if supply > 0 else 0
    mul = 12 * lam * (1 - lam) ** 2
    div = sqrt(n_holders + 2)  # Per Eq. 4
    return max_cap * mul / div

def simulate_shares(
    n_honest, k_sybil,
    honest_capital, sybil_capital,
    iters_per_week=7, weeks=60, max_cap=1.0
):
    """Simulate share evolution with iteration cap."""
    n_total = n_honest + k_sybil
    honest_bal = [honest_capital / n_honest] * n_honest
    sybil_bal = [sybil_capital / k_sybil] * k_sybil
    total_supply = honest_capital + sybil_capital

    trajectory = []
    for week in range(weeks):
        for _ in range(iters_per_week):
            honest_bal = [
                b + beta_cap(b, max_cap, total_supply, n_total)
                for b in honest_bal
            ]
            sybil_bal = [
                b + beta_cap(b, max_cap, total_supply, n_total)
                for b in sybil_bal
            ]
            total_supply = sum(honest_bal) + sum(sybil_bal)

        honest_share = sum(honest_bal) / total_supply
        sybil_share = sum(sybil_bal) / total_supply
        trajectory.append((week, honest_share, sybil_share))

    return trajectory
\end{lstlisting}

The simulation reveals that capital shares are \emph{persistent}, not convergent: the beta cap function preserves relative capital ratios because capacity gain is proportional to current $\lambda = B/S$.

\begin{table}[H]
\centering
\caption{Simulated share evolution (100 honest vs.\ 50 Sybil accounts)}
\label{tab:share-evolution}
\small
\begin{tabularx}{0.95\columnwidth}{@{}>{\centering\arraybackslash}X >{\centering\arraybackslash}X >{\centering\arraybackslash}X >{\centering\arraybackslash}X@{}}
\toprule
Scenario & Initial Share & Week 7 & Week 59 \\
\midrule
Equal shares & 50\%:50\% & 50.9\%:49.1\% & 51.4\%:48.6\% \\
Uneq. shares & 10.5\%:89.5\% & 11.2\%:88.8\% & 11.9\%:88.1\% \\
\bottomrule
\end{tabularx}
\end{table}

The results show that initial capital distribution strongly influences long-term shares: with a $9\times$ capital advantage, the advantaged party retains 88\% share after 60~weeks. Shares change by less than 2\% over 420~iterations, and the weekly cap is iteration-count bounded (7~per week), so burst timing is irrelevant. The protocol's Sybil resistance therefore derives from the $\sqrt{n{+}2}$ divisor and iteration caps rather than from equilibrium convergence---early capital advantages persist, and the defense operates through structural constraints.

\subsection{Limitations}

The simulation assumes static holder count (in practice, $n$ fluctuates), continuous cap gains (real protocols have discrete rounding), no holder floor dynamics (the bootstrap phase with $n_{\text{real}} < n_{\min}$ is unmodeled), and homogeneous iteration timing (real-world timing variations may introduce minor perturbations). Specifically, \texttt{Position.largeHolders()} returns $\max(\text{actual},\, \texttt{MIN\_HOLDERS\_ID})$; results in Tab.~\ref{tab:simulation-output} assume \texttt{MIN\_HOLDERS\_ID} $= 0$ (no floor). With the deployed floor, iteration counts at low $n$ are uniformly larger by the ratio $\sqrt{(\texttt{MIN\_HOLDERS\_ID} + 2)/(n + 2)}$.
The simulation validates two key properties from \S4.4 of~\cite{wp}: (1)~the $O(\sqrt{n})$ iteration scaling, and (2)~starting balance independence for single-account accumulation time. It also reveals that multi-account share dynamics do not converge to account-proportional equilibrium---a finding that informs the rate-limiting analysis in \S4.4 of~\cite{wp}.

\section{Cascade Simulation Implementation}
\label{apx:cascade-sim}

This appendix provides the cascade simulation implementation. Both scripts share a common \texttt{simulate} function and output semicolon-separated values.

\subsection{Multi-Scenario Comparison}

The following script compares lock effectiveness across five market liquidity scenarios at a fixed 20\% price shock:

\begin{lstlisting}[language=Python, caption={Multi-scenario cascade simulation}, label=lst:cascade-multi]
#!/usr/bin/env python3
from numpy import clip, random

def simulate(positions, price_drop, impact_coef):
    """Simulates liquidation cascade with market impact."""
    initial_supply = sum(p[0] for p in positions)
    price = 1.0 - price_drop
    total_liquidated = 0
    while True:
        underwater = [p for p in positions
            if p[0] > 0 and (p[0] * price) / p[1] < 1.0]
        if not underwater:
            break
        round_sold = 0
        for p in underwater:
            unlocked = p[0] * (1 - p[2])
            round_sold += unlocked
            total_liquidated += p[0]
            p[0] = 0
            p[1] = 0
        price -= impact_coef * round_sold
        if price <= 0:
            break
    return (total_liquidated / initial_supply * 100
            ) if initial_supply > 0 else 0

def gen_positions(seed, n, lock):
    """Generate positions with given lock fraction."""
    random.seed(seed)
    sizes = random.lognormal(mean=0, sigma=1, size=n)
    healths = clip(random.normal(1.5, 0.3, n), 1.0, 3.0)
    return [[s, s/h, lock] for s, h in zip(sizes, healths)]

if __name__ == "__main__":
    n, seed, drop = 1000, 42, 0.20
    scenarios = [("Liquid", 1e-6, 0.2),
                 ("Moderate", 6e-5, 10.1),
                 ("Strong", 1.5e-4, 25.2),
                 ("Severe", 3e-4, 50.5),
                 ("Extreme", 5e-4, 84.1)]
    print("scenario;depth;k;liq0;liq1;reduction")
    for name, k, depth in scenarios:
        liq0 = simulate(gen_positions(seed, n, 0.0), drop, k)
        liq1 = simulate(gen_positions(seed, n, 1.0), drop, k)
        reduction = (liq0-liq1)/liq0 * 100 if liq0 > 0 else 0
        print(f"{name};{depth:.1f};{k:.0e};{liq0:.1f};{liq1:.1f};{reduction:.1f}")
\end{lstlisting}

\subsection{Detailed Per-Scenario Analysis}

The following script generates detailed results for a specific scenario, varying both lock fraction and price shock:

\begin{lstlisting}[language=Python, caption={Detailed cascade simulation}, label=lst:cascade-detail]
#!/usr/bin/env python3
import sys
from numpy import clip, random

# simulate() and gen_positions() same as above (omitted)

SCENARIOS = [("Liquid", 1e-6, 0.2),
             ("Moderate", 6e-5, 10.1),
             ("Strong", 1.5e-4, 25.2),
             ("Severe", 3e-4, 50.5),
             ("Extreme", 5e-4, 84.1)]

if __name__ == "__main__":
    idx = int(sys.argv[1])  # 0-4
    name, k, depth = SCENARIOS[idx]
    n, seed = 1000, 42
    locks = [0.00, 0.25, 0.50, 0.75, 1.00]
    drops = [0.10, 0.15, 0.20, 0.25, 0.30]
    print("lock;drop;liquidated")
    for lock in locks:
        for drop in drops:
            result = simulate(gen_positions(seed, n, lock), drop, k)
            print(f"{lock:.2f};{drop:.2f};{result:.3f}")
\end{lstlisting}

\subsection{Algorithm}

The \texttt{simulate} function iterates a liquidation cascade: after applying the initial price drop, it identifies underwater positions ($H < 1$), liquidates them (computing unlocked collateral as supply $\times\,(1{-}\phi)$), applies linear market impact $p \leftarrow p - k \cdot V_{\text{sold}}$, and repeats until no underwater positions remain or price reaches zero. Here $H$ denotes the weighted health $H_{\text{ltv}} = (\text{supply} \cdot w_s)/(\text{borrow} \cdot w_b)$ that the contract checks at \texttt{Pool.sol:519} (\texttt{wnav\_supply < wnav\_borrow}). The simulation does not separately parametrise $(w_s, w_b)$; results are LTV-agnostic in this normalisation.

\subsection{Market Impact Model}

The linear price impact coefficient $k$ determines market depth ($1/k$ in supply units). The five scenarios range from \emph{Liquid} ($k = 10^{-6}$, pool is 0.2\% of depth, no cascade amplification) through \emph{Strong} ($k = 1.5 \times 10^{-4}$, 25\% of depth, 46\% reduction from locks) to \emph{Extreme} ($k = 5 \times 10^{-4}$, 84\% of depth, 80\% reduction). The locked fraction $\phi$ attenuates market impact by factor $(1{-}\phi)$, directly reducing cascade sell pressure.

\subsection{Limitations}

The simulation assumes complete (rather than partial $2^{-e}$) liquidation, a linear impact model (real markets exhibit convex slippage), greedy liquidators (real behavior may involve front-running or strategic timing), and a single collateral asset (cross-collateral pools may exhibit different dynamics).

\section{TWAP Oracle Simulations}
\label{apx:twap-sim}

This appendix presents simulation results for the log-space EMA smoothing and logarithmic spread scaling mechanisms. The oracle stores $\log_2(\text{mid})$ and $\log_2(1 + s_{\text{geo}})$, smoothed via EMA in log space. All simulations use 18-decimal precision ($1\ \text{unit} = 10^{18}$), matching the Solidity implementation.

\subsection{Decay Factor Visualization}

The EMA decay factor $\lambda = 0.5^{1/h}$ determines how quickly the TWAP responds to price changes, where $h$ is the half-life in refresh periods. As Figure~\ref{fig:app-decay} shows, $\lambda$ approaches 1.0 asymptotically with increasing half-life. The steep initial rise means that small changes in low $h$ values have large effects on responsiveness, while values above $h = 20$ yield diminishing returns. In practice, $\lambda$ values between 0.9 and 0.98 provide the useful operating range: below 0.9 the oracle becomes too reactive, while above 0.98 response times may be too slow for timely liquidations.

\begin{figure}[H]
\centering
\includegraphics[width=\columnwidth]{sims/sim_twap/plt-twap.decay.pdf}
\caption{Decay factor $\lambda$ by half-life. Higher half-life values produce decay factors closer to 1.0, making the TWAP more stable but slower to respond. The protocol default uses $h=12$ periods (\texttt{DECAY\_12HL} $= 0.5^{1/12} \approx 0.943874$, rounded to $\lambda \approx 0.944$ in subsequent prose).}
\label{fig:app-decay}
\end{figure}

\subsection{Price Shock Response}

Figures~\ref{fig:app-refresh-response} and~\ref{fig:app-refresh-decay} demonstrate how the TWAP mean responds to a sudden price shock (100 $\rightarrow$ 150) under different half-life configurations.

The half-life $h$ controls the trade-off between responsiveness and manipulation resistance. Shorter half-lives track market movements closely but are more vulnerable to manipulation; longer half-lives filter transient spikes but delay response to genuine shifts. The log-space EMA update $\overline{q}_{n+1} = \lambda \cdot \overline{q}_n + (1{-}\lambda) \cdot q_n$ ensures smooth convergence without overshooting, with the log transform dampening large deviations more aggressively than a linear EMA. As shown in Figure~\ref{fig:app-refresh-response}, even with HL=24 the TWAP reaches 90\% of the new price within approximately 55~periods.

\begin{figure}[H]
\centering
\includegraphics[width=\columnwidth]{sims/sim_twap/plt-twap.refresh-response.pdf}
\caption{TWAP response to a 50\% price shock at period 5. Shorter half-lives (HL=2) converge quickly but offer less manipulation resistance; longer half-lives (HL=24) provide robust smoothing but respond slowly.}
\label{fig:app-refresh-response}
\end{figure}

The memory decay $\lambda^n$ quantifies how quickly historical prices lose influence: after $k$ half-lives the residual weight is $0.5^k$ (12.5\% after 3, $<$1\% after 7). For $h = 12$, an attacker must sustain artificial prices for approximately 40~periods to shift the TWAP by 90\%.

\begin{figure}[H]
\centering
\includegraphics[width=\columnwidth]{sims/sim_twap/plt-twap.refresh-decay.pdf}
\caption{Memory decay showing the weight on pre-shock price over time. After HL periods, the weight on the old value drops to 50\%. Higher half-lives retain more memory of historical prices.}
\label{fig:app-refresh-decay}
\end{figure}

\begin{figure}[H]
\centering
\includegraphics[width=\columnwidth]{sims/sim_twap/plt-twap.liquidity.pdf}
\caption{Spread scaling by liquidity level. Illiquid pairs (high base spread) experience faster spread growth with trade size.}
\label{fig:app-liquidity}
\end{figure}

\subsection{Spread Scaling by Liquidity}

The effective spread depends on both position size and market liquidity (captured by the base spread $\sigma$). Figure~\ref{fig:app-liquidity} shows how different liquidity levels affect spread scaling.

The base spread $\sigma$ serves as a proxy for market depth. The formula $\mu = \log_2(2x + 2)$ with $x = n \cdot \sigma$ provides self-calibrating behavior: liquid markets (low $\sigma$) tolerate larger positions before significant spread widening, eliminating the need for per-pair parameter tuning while ensuring conservative valuations in thin markets.

\subsection{Complete Spread Analysis}

Figures~\ref{fig:app-spread-quotes}, \ref{fig:app-spread-percent}, and~\ref{fig:app-spread-multiplier} provide a comprehensive view of the spread scaling mechanism using a base spread of 2\%.

Together, these three views illustrate that while large positions incur wider spreads, the logarithmic scaling keeps them economically reasonable.
\begin{figure}[H]
\centering
\includegraphics[width=\columnwidth]{sims/sim_twap/plt-twap.spread-quotes.pdf}
\caption{Bid/ask quotes diverge from mid-value as trade size increases. The spread between bid and ask widens logarithmically, protecting against large position manipulation.}
\label{fig:app-spread-quotes}
\end{figure}

\begin{figure}[H]
\centering
\includegraphics[width=\columnwidth]{sims/sim_twap/plt-twap.spread-percent.pdf}
\caption{Effective spread grows logarithmically from the base spread (2.0\%). At 1 unit the spread is 2.0\%; at 1{,}000 units it reaches 10.8\%; at 1{,}000{,}000 units it reaches 30.6\%.}
\label{fig:app-spread-percent}
\end{figure}

\begin{figure}[H]
\centering
\includegraphics[width=\columnwidth]{sims/sim_twap/plt-twap.spread-multiplier.pdf}
\caption{The multiplier $\mu = \log_2(2x + 2)$ as a function of $x = n \cdot s$. This function provides bounded logarithmic growth: $\mu = 1$ at $x = 0$, growing to $\mu \approx 21$ at $x = 10^6$.}
\label{fig:app-spread-multiplier}
\end{figure}

\subsection{On-Chain Scenario Testing}
\label{subsec:oracle-scenarios}

To complement the preceding simulations, we conducted comprehensive on-chain scenario testing of the deployed oracle contract using Foundry. The test suite exercises 10~attack and stress scenarios across 6~token/liquidity configurations (60~total tests), measuring mid-price deviation and spread response under adversarial conditions. All tests use the production \texttt{DECAY\_12HL} ($\alpha \approx 0.944$, half-life $h{=}12$ refreshes) with 1-hour refresh intervals. Tables~\ref{tab:scenario-mid} and~\ref{tab:scenario-spread} summarize the results; all values are averages across the 6~configurations unless noted.

\subsubsection{Token/Liquidity Configurations}

The test matrix covers symmetric and asymmetric reserves across all supported decimal combinations:

\begin{table}[H]
\centering
\small
\caption{Oracle scenario test configurations}
\label{tab:scenario-configs}
\begin{tabularx}{0.95\columnwidth}{@{}c c X@{}}
\toprule
Decimals & Reserves & Profile \\
\midrule
18/18 & 100/100 & Symmetric standard \\
18/18 & 180k/50 & Asymmetric standard \\
6/6 & 100/100 & Symmetric stablecoin \\
6/6 & 180k/50 & Asymmetric stablecoin \\
18/6 & 100/100 & Mixed (18$\to$6) \\
6/18 & 100/100 & Mixed (6$\to$18) \\
\bottomrule
\end{tabularx}
\end{table}

\subsubsection{Mid-Price Manipulation Resistance}

Table~\ref{tab:scenario-mid} records the mid-price deviation $\Delta\%$ from baseline after each scenario, measuring how far an attacker (or legitimate market move) can shift the oracle's reported price.

\begin{table}[H]
\centering
\small
\caption{Mid-price deviation (\%) after each scenario}
\label{tab:scenario-mid}
\begin{tabularx}{0.95\columnwidth}{@{}l X r@{}}
\toprule
Scenario & Description & $\Delta\%$ \\
\midrule
S01 & Benign $\pm$5\% oscillations, 24 refreshes & 2.3\% \\
S02 & Instant 2$\times$ jump, 14 sustained refreshes & 22.5\% \\
S03 & Gradual 2$\times$ drift over 12 refreshes & 9.1\% \\
S04 & 2$\times$ spike, immediate return to baseline & 3.7\% \\
S05 & 2$\times$ spike, gradual return over 12 refreshes & 14.9\% \\
S06 & 10$\times$ liquidity increase (same ratio) & ${<}0.01$\% \\
S07 & Gradual 10$\times$ liquidity increase & ${<}0.01$\% \\
S08 & 90\% liquidity drain (same ratio) & ${<}0.06$\% \\
S09 & Gradual 90\% liquidity drain & ${<}0.01$\% \\
S10a & 3$\times$ manipulation between refreshes & \textbf{0\%} \\
S10b & 3$\times$ manipulation at single refresh & \textbf{0\%} \\
\bottomrule
\end{tabularx}
\end{table}

\paragraph{Key observations.}
The delayed EMA provides a \emph{two-tick immunity window}: the first two refreshes after any price spike produce exactly 0\% mid-price impact (S02, ticks~4--5), because the spike enters the \texttt{last} slot on the first refresh and only propagates into \texttt{mean} on the third (${\sim}4\%$ absorption per tick). After 14~sustained refreshes, only 22.5\% of a $2\times$ move has been absorbed, consistent with the theoretical bound from \S4.6 of~\cite{wp}. Inter-refresh manipulation (S10a) has exactly zero effect, as does single-tick manipulation at refresh time (S10b). Liquidity changes at constant price ratio (S06--S09) produce negligible mid-price deviation ($<$0.06\%), confirming that the log-space mid tracks price independently of pool depth.

\subsubsection{Spread Auto-Widening Response}

Table~\ref{tab:scenario-spread} records the spread percentage for a 1-unit query before and after liquidity changes. The spread correctly reflects pool depth without manual parameterization.

\begin{table}[H]
\centering
\small
\caption{Spread response to liquidity changes (1-unit query)}
\label{tab:scenario-spread}
\begin{tabularx}{0.95\columnwidth}{@{}l X r r@{}}
\toprule
Scenario & Description & Before & After \\
\midrule
S06 & Sudden 10$\times$ liquidity & 2.53\% & 2.37\% \\
S07 & Gradual 10$\times$ liquidity & 2.53\% & 1.98\% \\
S08 & Sudden 90\% drain & 2.06\% & 3.06\% \\
S09 & Gradual 90\% drain & 2.06\% & 2.52\% \\
\bottomrule
\end{tabularx}
\end{table}

EMA smoothing of the log-space spread produces an asymmetry between sudden and gradual changes: sudden liquidity addition yields $\approx$2.37\% after 3~ticks (partial absorption), while gradual addition over 12~ticks converges to $\approx$1.98\%. Sudden removal produces ${\sim}50\%$ spread widening (2.06\%\,$\to$\,3.06\%), providing immediate protection against thin-pool manipulation.

\subsubsection{Decimal Invariance}

All scenarios produce identical behavior (to within rounding) across the 6~token configurations. The sole exception is the (6/6, 180k/50) configuration, where integer truncation in 6-decimal arithmetic produces slightly coarser values (e.g., S04 residual of 3.99\% vs.\ 3.74\% elsewhere). This confirms that the log-space design is fundamentally decimal-agnostic.

\subsubsection{Strengths and Weaknesses}

\paragraph{Strengths.}
\begin{enumerate}[leftmargin=*]
\item \emph{Two-tick delay buffer}: The first two refreshes after a spike show 0\% mid-price impact, providing strong flash manipulation resistance.
\item \emph{Inter-refresh immunity}: Price changes between refresh windows have exactly zero effect on the oracle.
\item \emph{Self-calibrating spread}: The spread correctly widens and tightens with pool depth changes while keeping the mid stable, without per-pair configuration.
\item \emph{Decimal invariance}: Behavior is consistent across all token decimal combinations (6, 18, and mixed).
\end{enumerate}

\paragraph{Weaknesses (inherent to DECAY\_12HL).}
\begin{enumerate}[leftmargin=*]
\item \emph{Slow convergence}: After a legitimate $2\times$ price move, only ${\sim}22.5\%$ is absorbed after 14~hours, potentially enabling arbitrage against stale prices.
\item \emph{EMA reversal lag}: After a $2\times$ spike and full return to baseline, ${\sim}15\%$ residual remains, as the EMA retains memory of past spikes beyond their market relevance.
\item \emph{Slow drift undertracking}: A gradual $2\times$ move over 12~hours shows only 9.1\% oracle movement, creating a persistent gap between reported and actual price.
\end{enumerate}

These weaknesses are the intentional cost of heavy smoothing ($\alpha \approx 0.944$). The design prioritizes manipulation resistance over responsiveness---appropriate for a lending protocol where the cost of oracle manipulation (bad debt) exceeds the cost of lagging genuine price moves (delayed liquidations). The hourly refresh cadence further limits the attack surface, as an attacker cannot increase the sampling rate to accelerate EMA convergence toward a manipulated price.

% ==============================================================================
% Appendix G: Bad Debt Risk Quantification
% ==============================================================================
\section{Bad Debt Risk Quantification}
\label{apx:bad-debt}

\noindent\textit{This appendix was originally published as a standalone companion document and is reproduced here for completeness.}

\subsection{Introduction}
\label{sec:bd-intro}

XPower Banq's log-space TWAP oracle provides robust manipulation resistance---approximately 40~hours of sustained artificial pricing are required to achieve 90\% deviation~\cite{mackinga2022}. However, the same smoothing that resists manipulation delays response to genuine market crashes. After an instantaneous price crash, only ${\sim}50\%$ of the log-space deviation is absorbed after 13~hourly refreshes (the 12-hour half-life). Combined with the 1-hour rate limit between oracle updates, the protocol can be 2+~hours blind during a crash.

This appendix provides a complete risk quantification framework:
\begin{enumerate}[leftmargin=*]
\item \textbf{Oracle Lag Model} (\S\ref{sec:bd-oracle-lag}): Formal characterization of the EMA convergence delay and the resulting phantom-healthy window.
\item \textbf{Monte Carlo Simulation} (\S\ref{sec:bd-montecarlo}): Jump-diffusion price paths with ETH-calibrated parameters, swept across LTV and oracle decay configurations.
\item \textbf{Analytical Bound} (\S\ref{sec:bd-bound}): Closed-form upper bound on bad debt as a function of crash magnitude, oracle decay, and LTV.
\item \textbf{Safe Operating Region} (\S\ref{sec:bd-safety}): Parameter configurations satisfying the $<$5\% TVL constraint.
\end{enumerate}

All simulations use fixed random seeds for reproducibility.

% ==============================================================================
% 2. Oracle Lag Model
% ==============================================================================
\subsection{Oracle Lag Model}
\label{sec:bd-oracle-lag}

\subsubsection{EMA Update Semantics}
\label{subsec:bd-ema}

The TWAP oracle (\texttt{TWAP.sol}) performs an EMA update on each refresh. Critically, the update blends the stored mean with the \emph{previous} observation (\texttt{last}), not the current one (\texttt{next}):
\begin{align}
\overline{m}_n &= \alpha \cdot \overline{m}_{n-1} + (1{-}\alpha) \cdot \ell_{n-1} \label{eq:ema-update} \\
\ell_n &= q_n \label{eq:last-update}
\end{align}
where $\overline{m}_n$ is the smoothed log-price, $\ell_n$ is the stored last observation, and $q_n$ is the current market price (in $\log_2$ space). This creates a \emph{one-refresh delay}: a new price enters the mean only on the \emph{following} refresh.

\subsubsection{Step-Crash Convergence}
\label{subsec:bd-step-crash}

After an instantaneous crash of fraction $\delta$ at time $t_0$ (price drops from $p_0$ to $p_0(1{-}\delta)$), the oracle-reported price after $n$ refreshes follows:
\begin{equation}
\label{eq:oracle-trajectory}
\hat{p}(n) = p_0 \cdot 2^{\overline{m}_n}
\end{equation}
where the log-space EMA converges geometrically toward the post-crash price. The fraction of the log-space deviation absorbed by the oracle at refresh $n$ is:
\begin{equation}
\label{eq:absorption}
A(n) = 1 - \alpha^{n-1} \quad (n \geq 2), \qquad A(0) = A(1) = 0
\end{equation}

Due to the one-refresh delay, the oracle is completely blind for the first refresh after a crash. This log-space absorption is independent of crash magnitude $\delta$---all crashes converge at the same geometric rate $(1{-}\alpha)$ per refresh. The oracle-reported price in linear space is $\hat{p}(n) = p_0 \cdot (1{-}\delta)^{A(n)}$ (\texttt{mean.mid} is updated independently of \texttt{mean.rel}), so linear-space absorption varies with $\delta$. Figure~\ref{fig:oracle-convergence} shows the convergence trajectories for various crash magnitudes, and Table~\ref{tab:absorption} shows log-space absorption percentages for the default decay $\alpha = 0.944$ (12-hour half-life).

\begin{figure}[htbp]
\centering
\includegraphics[width=\columnwidth]{sims/sim_debt/output/01b_oracle_convergence.pdf}
\caption{Oracle convergence after step crashes of various magnitudes ($\alpha = 0.944$, 1\,h refresh). The dashed line marks the 12-hour half-life. All crash magnitudes converge at the same geometric rate in log-space, reaching 50\% absorption near $n = 13$ refreshes (accounting for the one-refresh delay).}
\label{fig:oracle-convergence}
\end{figure}

\begin{table}[htbp]
\centering
\caption{Oracle absorption $A(n)$ after $n$ hourly refreshes ($\alpha = 0.944$)}
\label{tab:absorption}
\small
\begin{tabularx}{0.95\columnwidth}{@{}r *{5}{>{\centering\arraybackslash}X}@{}}
\toprule
$n$ & 20\% & 30\% & 40\% & 50\% & 70\% \\
\midrule
1  & 0.0\% & 0.0\% & 0.0\% & 0.0\% & 0.0\% \\
2  & 5.6\% & 5.6\% & 5.6\% & 5.6\% & 5.6\% \\
4  & 15.9\% & 15.9\% & 15.9\% & 15.9\% & 15.9\% \\
8  & 33.2\% & 33.2\% & 33.2\% & 33.2\% & 33.2\% \\
12 & 46.9\% & 46.9\% & 46.9\% & 46.9\% & 46.9\% \\
24 & 73.4\% & 73.4\% & 73.4\% & 73.4\% & 73.4\% \\
48 & 93.3\% & 93.3\% & 93.3\% & 93.3\% & 93.3\% \\
72 & 98.3\% & 98.3\% & 98.3\% & 98.3\% & 98.3\% \\
\bottomrule
\end{tabularx}
\vspace{0.3em}
\par\raggedright\scriptsize Log-space absorption $A(n) = 1 - \alpha^{n-1}$ is independent of crash magnitude; linear-space absorption varies with $\delta$.
\end{table}

\begin{figure}[htbp]
\centering
\includegraphics[width=\columnwidth]{sims/sim_debt/output/01_oracle_lag_heatmap.pdf}
\caption{Oracle deviation from true price (\%) as a function of crash magnitude and hours elapsed. The heatmap shows the full convergence surface; larger crashes produce higher absolute deviations but the same geometric decay rate.}
\label{fig:oracle-lag-heatmap}
\end{figure}

\subsubsection{Phantom-Healthy Windows}
\label{subsec:bd-phantom}

\begin{definition}[Phantom-Healthy Position]
\label{def:phantom}
A position is \emph{phantom-healthy} at time $t$ when:
\begin{align}
H_{\text{oracle}}(t) &= H_0 \cdot \frac{\hat{p}(t)}{p_0} \geq 1 \label{eq:h-oracle} \\
H_{\text{true}}(t) &= H_0 \cdot (1 - \delta) < 1 \label{eq:h-true}
\end{align}
where $H_0$ is the initial health factor.
\end{definition}

For a position to be truly underwater, the crash must satisfy $\delta > 1 - 1/H_0$. The oracle triggers liquidation when $\hat{p}(n)/p_0 < 1/H_0$, i.e., the oracle price ratio drops below the inverse of the initial health.

\begin{definition}[Phantom-Healthy Window]
The \emph{phantom-healthy window} $W$ is the number of oracle refreshes during which a position remains phantom-healthy:
\begin{equation}
W = \min\left\{n \geq 0 \;\middle|\; \frac{\hat{p}(n)}{p_0} < \frac{1}{H_0}\right\}
\end{equation}
\end{definition}

Table~\ref{tab:phantom} shows phantom-healthy windows for representative (crash, initial health) pairs with $\alpha = 0.944$.

\begin{table}[htbp]
\centering
\caption{Phantom-healthy window $W$ (hours) by crash magnitude and initial health, $\alpha = 0.944$}
\label{tab:phantom}
\small
\begin{tabularx}{0.95\columnwidth}{@{}r r *{4}{>{\centering\arraybackslash}X}@{}}
\toprule
& & \multicolumn{4}{c}{Initial Health $H_0$} \\
\cmidrule(l){3-6}
Crash & & 1.1 & 1.3 & 1.5 & 2.0 \\
\midrule
20\% & $W$ & 11 & --- & --- & --- \\
30\% & $W$ & 7 & 25 & --- & --- \\
40\% & $W$ & 5 & 14 & 29 & --- \\
50\% & $W$ & 4 & 10 & 17 & --- \\
70\% & $W$ & 3 & 6 & 9 & 16 \\
\bottomrule
\end{tabularx}
\vspace{0.3em}
\par\raggedright\scriptsize ``---'' indicates the position remains solvent (not underwater). Values verified against step-by-step OracleModel simulation (match within $\pm 1$ refresh).
\end{table}

\subsubsection{Worst-Case Blind Time}
\label{subsec:bd-blind}

The worst case occurs when a crash happens immediately after a refresh, adding one full refresh interval of blindness. The total blind time is:
\begin{equation}
\label{eq:blind-time}
t_{\text{blind}} = \Delta t_{\text{refresh}} + W \cdot \Delta t_{\text{refresh}}
\end{equation}

For the conservative-mode floor (33\% LTV) configuration with positions initialized near $H_0 = 1.5$:
\begin{itemize}[leftmargin=*]
\item \textbf{40\% crash}: $W = 29$ refreshes $\Rightarrow$ $t_{\text{blind}} = 30$~hours
\item \textbf{50\% crash}: $W = 17$ refreshes $\Rightarrow$ $t_{\text{blind}} = 18$~hours
\item \textbf{70\% crash}: $W = 9$ refreshes $\Rightarrow$ $t_{\text{blind}} = 10$~hours
\end{itemize}

\noindent Larger crashes produce shorter phantom windows because the true price diverges more rapidly from the oracle threshold. Positions with higher initial health (further from the liquidation boundary) require more refreshes before the oracle detects their distress.

% ==============================================================================
% 3. Monte Carlo Simulation
% ==============================================================================
\subsection{Monte Carlo Simulation}
\label{sec:bd-montecarlo}

\subsubsection{Price Process}
\label{subsec:bd-price}

Collateral prices follow a Merton jump-diffusion model~\cite{merton1976}:
\begin{equation}
\label{eq:jump-diffusion}
\frac{dp}{p} = \mu\,dt + \sigma\,dW + J\,dN
\end{equation}
where $W$ is a Wiener process, $N$ is a Poisson process with intensity $\lambda_J$, and $J \sim \mathcal{N}(\mu_J, \sigma_J^2)$ are i.i.d.\ jump magnitudes. Parameters are calibrated to ETH historical data (Table~\ref{tab:jd-params}).

\begin{table}[htbp]
\centering
\caption{Jump-diffusion parameters (ETH calibration)}
\label{tab:jd-params}
\small
\begin{tabularx}{0.85\columnwidth}{@{}>{\raggedright\arraybackslash}X r@{}}
\toprule
Parameter & Value \\
\midrule
Annual volatility $\sigma$ & 90\% \\
Jump intensity $\lambda_J$ (jumps/year) & 6.0 \\
Mean jump $\mu_J$ & $-15\%$ \\
Jump volatility $\sigma_J$ & 10\% \\
Drift $\mu$ & compensated \\
\bottomrule
\end{tabularx}
\end{table}

\subsubsection{Simulation Design}
\label{subsec:bd-sim_design}

\textbf{Resolution.} Paths are generated at hourly resolution (8{,}760 steps per year), matching the oracle refresh cadence. This is a deliberate trade-off: minute-level resolution would require $100\text{K} \times 525\text{K} = 52.5$B data points (${\sim}210$~GB at float32), while hourly resolution requires only 875M points (${\sim}3.5$~GB). Since the oracle can only update hourly, intra-hour price movements do not affect oracle-triggered liquidations.

\textbf{Pool.} Each path initializes 1{,}000 synthetic positions with health factors $H_0 \sim \mathcal{N}(1.5, 0.3)$ truncated at $[1.0, 3.0]$ and uniform unit borrow size.

\textbf{Oracle.} The log-space EMA replicates \texttt{TWAP.sol} semantics exactly, including the one-refresh delay. The oracle processes hourly price snapshots and produces oracle-price-ratio time series.

\textbf{Refresh.} The 1-hour cadence assumed throughout this appendix is the deployed value of the governance parameter \texttt{LIMIT\_ID} (set in the deploy script and gated by the \texttt{delayed()} modifier on \texttt{Oracle.refresh()}); it is not a hard-coded constant. The conservative-mode parameter set and the phantom-window math (Eq.~\ref{eq:blind-time}, Tab.~\ref{tab:phantom}) scale linearly with $\Delta t_{\text{refresh}}$ and are contingent on this cadence not being shortened materially.

\textbf{Liquidation.} Two models are implemented: (1)~\emph{full liquidation} via vectorized threshold crossing---for each position, pre-compute the critical oracle price ratio $p_{\text{crit}} = 1/H_0$ and detect the first crossing with \texttt{np.argmax}; (2)~\emph{partial liquidation} that transfers fraction $2^{-e}$ per step (matching \texttt{Pool.sol}'s right-shift by \texttt{partial\_exp}; $e = 1$ gives the default 50\%-per-step), simulated step-by-step with optional dimensionless liquidation-recovery haircut $\kappa$ (distinct from the cascade-simulation depth coefficient $k$ of \S\ref{apx:cascade-sim}).

\textbf{Parameters.} Default run: 1{,}000 paths, LTV = 1/3 (the conservative-mode governance floor), $\alpha = 0.944$, 1-hour refresh, seed = 42. The full LTV sweep---including the shipped default of 2/3 (66.67\%)---appears in Table~\ref{tab:mc-ltv}; per-path detailed plots are rendered at the 1/3 floor as a worst-case visual demonstration.

\subsubsection{Results}
\label{subsec:bd-mc-results}

Table~\ref{tab:mc-ltv} presents Monte Carlo results across LTV configurations. The shipped default of 66.67\% LTV sits at the 67\% row, with $\mathbb{E}[\text{BD}] = 0.00\%$, CVaR(99\%) = 0.02\%, and 1.3\% of paths producing any bad debt---bounded but non-zero risk under tail-event jump-diffusion paths. The conservative-mode floor of 33\% LTV produces zero bad debt across all paths and metrics: the 200\% over-collateralization buffer ensures the oracle triggers liquidation (at $p_{\text{crit}} = 1/H_0$) well before any shortfall accumulates. Bad debt grows non-linearly above 67\%, reaching $\mathbb{E}[\text{BD}] = 0.01\%$ and CVaR(99\%) = 0.69\% at the 75\% upper sweep point.

\begin{table}[htbp]
\centering
\caption{Monte Carlo bad debt metrics by LTV ($\alpha = 0.944$, $N = 1{,}000$ paths)}
\label{tab:mc-ltv}
\small
\begin{tabularx}{0.95\columnwidth}{@{}r *{4}{>{\centering\arraybackslash}X}@{}}
\toprule
LTV & $\mathbb{E}[\text{BD}]$ & VaR(99\%) & CVaR(99\%) & Nonzero \\
\midrule
33\% & 0.00\% & 0.00\% & 0.00\% & 0.0\% \\
50\% & 0.00\% & 0.00\% & 0.00\% & 0.0\% \\
67\% & 0.00\% & 0.00\% & 0.02\% & 1.3\% \\
75\% & 0.01\% & 0.32\% & 0.69\% & 10.7\% \\
\bottomrule
\end{tabularx}
\vspace{0.3em}
\par\raggedright\scriptsize All values as \% of TVL per year. ``Nonzero'' is the fraction of paths producing any bad debt. 95\% CI on $\mathbb{E}[\text{BD}]$ is $< \pm 0.10\%$ at 1{,}000 paths.
\end{table}

Table~\ref{tab:mc-alpha} shows the effect of oracle decay rate. At the 33\% LTV floor, all decay rates produce zero bad debt---the over-collateralization margin dominates. Oracle decay becomes relevant only at higher LTV (including the shipped 66.67\% default) where the liquidation-to-bad-debt gap is narrower.

\begin{table}[htbp]
\centering
\caption{Monte Carlo bad debt metrics by oracle decay ($\text{LTV} = 33\%$, $N = 1{,}000$ paths)}
\label{tab:mc-alpha}
\small
\begin{tabularx}{0.95\columnwidth}{@{}r r *{3}{>{\centering\arraybackslash}X}@{}}
\toprule
$\alpha$ & HL & $\mathbb{E}[\text{BD}]$ & VaR(99\%) & CVaR(99\%) \\
\midrule
0.891 & 6\,h & 0.00\% & 0.00\% & 0.00\% \\
0.944 & 12\,h & 0.00\% & 0.00\% & 0.00\% \\
0.972 & 24\,h & 0.00\% & 0.00\% & 0.00\% \\
\bottomrule
\end{tabularx}
\vspace{0.3em}
\par\raggedright\scriptsize HL = half-life in hours. At the 33\% LTV floor, the 200\% over-collateralization absorbs all oracle staleness regardless of decay rate; results at the shipped 66.67\% default are sensitive to $\alpha$ (see Tab.~\ref{tab:mc-ltv}).
\end{table}

\subsubsection{Bad Debt Distribution}
\label{subsec:bd-distribution}

Figure~\ref{fig:mc-histogram} shows the full bad debt distribution for the conservative-mode (33\% LTV) configuration. All 1{,}000 simulated paths produce exactly zero bad debt, confirming that the 200\% over-collateralization at the 33\% LTV floor fully absorbs oracle staleness. The oracle triggers liquidation at $p_{\text{crit}} = 1/H_0 \approx 0.67$ (33\% price decline), well before the bad-debt threshold at $p = \text{LTV}/H_0 \approx 0.22$ (78\% decline). At the shipped default of 66.67\% LTV the bad-debt threshold rises to $p \approx 0.44$ (56\% decline), reducing the gap and accounting for the 1.3\% of nonzero-BD paths in Table~\ref{tab:mc-ltv}.

\begin{figure}[htbp]
\centering
\includegraphics[width=\columnwidth]{sims/sim_debt/output/04_mc_histogram_quick.pdf}
\caption{Bad debt distribution at the conservative-mode floor (LTV = 33\%, $\alpha = 0.944$). All paths produce zero bad debt: the 200\% over-collateralization ensures liquidation triggers well before any shortfall accumulates. The shipped default at 66.67\% LTV produces bounded but non-zero bad debt; see Tab.~\ref{tab:mc-ltv}.}
\label{fig:mc-histogram}
\end{figure}

\subsubsection{Drawdown--Bad Debt Relationship}
\label{subsec:bd-drawdown}

Figure~\ref{fig:drawdown-bd} confirms that at the 33\% LTV floor, even paths with extreme drawdowns ($>$90\%) produce zero bad debt. Bad debt requires the \emph{true} price to drop below $\text{LTV}/H_0 \approx 0.22$ \emph{and} the oracle to fail to trigger liquidation in time---the corrected liquidation threshold ($1/H_0 \approx 0.67$) catches positions long before this point.

\begin{figure}[htbp]
\centering
\includegraphics[width=\columnwidth]{sims/sim_debt/output/04b_drawdown_vs_bd_quick.pdf}
\caption{Maximum drawdown vs.\ bad debt per simulated path at the conservative-mode floor (LTV = 33\%, $\alpha = 0.944$). All paths produce zero bad debt regardless of drawdown severity, confirming the analytical bound.}
\label{fig:drawdown-bd}
\end{figure}

\subsubsection{Liquidation Delay Distribution}
\label{subsec:bd-liq-delay}

For paths that produce liquidations, the delay between a position becoming truly underwater ($H_{\text{true}} < 1$) and the oracle triggering liquidation ($H_{\text{oracle}} < 1$) has the following characteristics (baseline configuration, 1{,}000 paths):

\begin{itemize}[leftmargin=*]
\item \textbf{Median delay}: 30~hours
\item \textbf{Mean delay}: $>$218~hours (skewed by long tails)
\item \textbf{99th percentile}: ${\sim}$3{,}800~hours
\end{itemize}

The heavy right tail (Figure~\ref{fig:liq-delay}) reflects paths where positions hover near the liquidation boundary during gradual declines, maintaining phantom-healthy status for extended periods.

\begin{figure}[htbp]
\centering
\includegraphics[width=\columnwidth]{sims/sim_debt/output/04c_time_to_liquidation_quick.pdf}
\caption{Distribution of oracle-induced liquidation delay (hours from true underwater to oracle-triggered liquidation). The distribution is extremely heavy-tailed: median delay is 30\,h but the mean exceeds 218\,h due to paths where positions hover near the liquidation boundary.}
\label{fig:liq-delay}
\end{figure}

% ==============================================================================
% 4. Analytical Bound Derivation
% ==============================================================================
\subsection{Analytical Bound}
\label{sec:bd-bound}

\subsubsection{Instant Liquidation Baseline}
\label{subsec:bd-instant}

As a baseline, consider bad debt under instant liquidation (no oracle delay). A position with initial health $H_0$ experiencing crash fraction $\delta$ has:
\begin{equation}
\text{BD}_{\text{inst}} = \text{borrow} \cdot \max\!\left(0,\; 1 - \frac{H_0(1{-}\delta)}{\text{LTV}}\right)
\end{equation}

Positions are underwater when $H_0 < 1/(1{-}\delta)$, and produce bad debt when $H_0(1{-}\delta) < \text{LTV}$ (shortfall exceeds collateral value at true price). Integrating over the health distribution $H_0 \sim \text{TruncNorm}(1.5, 0.3, [1.0, 3.0])$ and weighting by TVL contribution:
\begin{equation}
\label{eq:bd-instant}
\frac{\mathbb{E}[\text{BD}_{\text{inst}}]}{\text{TVL}} = \frac{\int_{1}^{H_{\text{crit}}} \max\!\left(0,\; 1 - \frac{H(1{-}\delta)}{\text{LTV}}\right) H \cdot f(H)\,dH}{\mathbb{E}[H]}
\end{equation}
where $H_{\text{crit}} = 1/(1{-}\delta)$ and $f(H)$ is the truncated normal PDF.

\subsubsection{Oracle Delay Penalty}
\label{subsec:bd-delay}

The additional bad debt from oracle staleness arises when price continues declining during the phantom-healthy window. For a position with phantom window $W$ refreshes and hourly volatility $\sigma_h = \sigma/\sqrt{8760}$, the expected additional decline is:
\begin{equation}
\label{eq:additional-decline}
\Delta_W = W \cdot \sigma_h \approx W \cdot \frac{0.90}{\sqrt{8760}} \approx 0.0096 \cdot W
\end{equation}

\subsubsection{Conservative Upper Bound}
\label{subsec:bd-conservative}

\begin{theorem}[Bad Debt Upper Bound]
\label{thm:bd-bound}
For a step crash of fraction $\delta$, oracle decay $\alpha$, and effective LTV, the maximum bad debt as a fraction of TVL is bounded by:
\begin{equation}
\label{eq:bd-bound}
\text{BD}_{\max}(\delta, \alpha, \text{LTV}) \leq \text{BD}_{\text{inst}} \cdot \left(1 + \frac{W_{\max} \cdot \sigma_h}{\delta}\right)
\end{equation}
where $W_{\max}$ is the phantom window for the marginal underwater position (health just below $1/(1{-}\delta)$) and $\sigma_h \approx 0.96\%$/hour is the ETH hourly volatility.
\end{theorem}

\begin{proof}
The bound follows from three observations: (1)~all positions that are underwater under instant liquidation are also underwater under delayed liquidation, establishing $\text{BD}_{\text{inst}}$ as a lower bound; (2)~positions with longer phantom windows experience additional price movement of order $W \cdot \sigma_h$ during the delay; (3)~the marginal underwater position (highest $H_0$ still underwater) has the longest phantom window $W_{\max}$.

The multiplicative penalty $W_{\max} \cdot \sigma_h / \delta$ represents the worst-case fractional increase in bad debt due to continued decline during the phantom window, normalized by the initial crash magnitude. This is conservative because: (a)~it assumes all positions experience $W_{\max}$ rather than their individual (shorter) windows, and (b)~it assumes continued decline at full volatility rather than zero-drift random walk.
\end{proof}

Table~\ref{tab:bound} presents the bound for selected parameter combinations.

\begin{table}[htbp]
\centering
\caption{Analytical bad debt bound (\% of TVL)}
\label{tab:bound}
\small
\begin{tabularx}{0.95\columnwidth}{@{}r r >{\centering\arraybackslash}X >{\centering\arraybackslash}X >{\centering\arraybackslash}X@{}}
\toprule
Crash & LTV & BD\textsubscript{inst} & $\alpha{=}0.944$ & $\alpha{=}0.891$ \\
\midrule
30\% & 33\% & 0.00 & 0.00 & 0.00 \\
30\% & 75\% & 0.06 & 0.20 & 0.13 \\
40\% & 67\% & 0.16 & 0.46 & 0.32 \\
40\% & 75\% & 0.97 & 2.79 & 1.90 \\
50\% & 33\% & 0.00 & 0.00 & 0.00 \\
50\% & 50\% & 0.00 & 0.00 & 0.00 \\
50\% & 67\% & 1.90 & 5.08 & 3.51 \\
50\% & 75\% & 4.98 & 13.32 & 9.20 \\
70\% & 33\% & 0.16 & 0.26 & 0.21 \\
70\% & 50\% & 9.67 & 15.52 & 12.73 \\
70\% & 67\% & 29.00 & 46.52 & 38.15 \\
\bottomrule
\end{tabularx}
\vspace{0.3em}
\par\raggedright\scriptsize BD\textsubscript{inst} is the baseline under instant liquidation. Bounds are computed with ETH hourly volatility $\sigma_h \approx 0.96\%$.
\end{table}

% ==============================================================================
% 5. Safe Operating Region
% ==============================================================================
\subsection{Safe Operating Region}
\label{sec:bd-safety}

\subsubsection{Safety Criterion}
\label{subsec:bd-criterion}

We define a parameter configuration $(\alpha, \text{LTV})$ as \emph{safe} if the analytical bad debt bound satisfies $\text{BD}_{\max} < 5\%$ of TVL for crash magnitudes up to 50\%. This corresponds to Acceptance Criterion~4 from the specification.

\subsubsection{Safe Configurations}
\label{subsec:bd-safe-configs}

Table~\ref{tab:safe} lists all configurations satisfying the safety criterion.

\begin{table}[htbp]
\centering
\caption{Safe parameter configurations (BD bound $< 5\%$ for 50\% crash)}
\label{tab:safe}
\small
\begin{tabularx}{0.95\columnwidth}{@{}r r >{\centering\arraybackslash}X r@{}}
\toprule
$\alpha$ & Half-Life & LTV & BD Bound \\
\midrule
0.891 & 6\,h & 33\% & 0.00\% \\
0.891 & 6\,h & 50\% & 0.00\% \\
0.891 & 6\,h & 67\% & 3.51\% \\
0.944 & 12\,h & 33\% & 0.00\% \\
0.944 & 12\,h & 50\% & 0.00\% \\
0.972 & 24\,h & 33\% & 0.00\% \\
0.972 & 24\,h & 50\% & 0.00\% \\
\bottomrule
\end{tabularx}
\vspace{0.3em}
\par\raggedright\scriptsize The conservative-mode floor ($\alpha = 0.944$, LTV = 33\%) produces zero bounded bad debt for any crash up to 50\%. The shipped default ($\alpha = 0.944$, LTV = 66.67\%) sits between the 50\% and 67\% rows and is governable down to 33\% in one lethargic cycle (halving $w_s$ from 170 to 85).
\end{table}

\begin{figure}[htbp]
\centering
\includegraphics[width=\columnwidth]{sims/sim_debt/output/08_safe_region.pdf}
\caption{Safe operating region in $(\alpha, \text{LTV}, \text{crash})$ space. The surface separates configurations where $\text{BD}_{\max} < 5\%$ (below) from those exceeding the threshold (above).}
\label{fig:safe-region}
\end{figure}

\subsubsection{Key Observations}
\label{subsec:bd-observations}

\begin{enumerate}[leftmargin=*]
\item \textbf{The 33\% LTV floor is zero-bad-debt for $\leq$50\% crashes.} At the conservative-mode floor, positions require a $>$50\% crash before any become underwater ($H_{\text{crit}} = 1/(1{-}\delta) > H_{\min} = 1.0$ requires $\delta > 0$, but $H_0 \cdot (1{-}\delta) < \text{LTV}$ requires $\delta > 1 - \text{LTV}/H_0$). With 200\% over-collateralization, even a 50\% crash leaves the average position ($H_0 = 1.5$) with $H_{\text{true}} = 0.75 > \text{LTV} = 0.33$: still solvent.

\item \textbf{LTV is the dominant risk parameter.} Bad debt is zero for LTV $\leq 50\%$, bounded but non-zero at the shipped 66.67\% default ($\mathbb{E}[\text{BD}] = 0.00\%$, CVaR(99\%) = 0.02\%), and grows non-linearly above, reaching $\mathbb{E}[\text{BD}] = 0.01\%$ and CVaR(99\%) = 0.69\% at the 75\% upper sweep point (Table~\ref{tab:mc-ltv}). The oracle decay rate has no observable effect at the 33\% LTV floor (Table~\ref{tab:mc-alpha}) but interacts non-linearly with LTV at the shipped default.

\item \textbf{The 50\% buffer absorbs oracle staleness within tolerance at the shipped default.} The oracle triggers liquidation at $p_{\text{crit}} = 1/H_0 \approx 0.67$ (33\% price decline). At the shipped 66.67\% LTV, bad debt requires the true price to reach $\text{LTV}/H_0 \approx 0.44$ (56\% decline)---a 23-percentage-point gap. At the conservative 33\% floor, the gap widens to 45 percentage points and oracle delay is rendered harmless.

\item \textbf{Tail risk is bounded at the shipped default and negligible at the floor.} At the shipped 66.67\% LTV, CVaR(99\%) = 0.02\%; at the conservative 33\% floor, all risk metrics (E[BD], VaR, CVaR) are zero. At the 75\% upper sweep point, CVaR(99\%) reaches 0.69\% of TVL.
\end{enumerate}

% ==============================================================================
% 6. Sensitivity Analysis
% ==============================================================================
\subsection{Sensitivity Analysis}
\label{sec:bd-sensitivity}

\subsubsection{Parameter Rankings}
\label{subsec:bd-tornado}

The following parameters are ranked by their impact on expected bad debt (based on tornado diagram analysis over the MC sweep):

\begin{enumerate}[leftmargin=*]
\item \textbf{LTV (weight ratio)} --- governance-controlled, dominant effect. Moving from the 33\% floor through the shipped 66.67\% default to the 75\% upper sweep point increases $\mathbb{E}[\text{BD}]$ from 0.00\% to 0.01\% and CVaR(99\%) from 0.00\% (at floor) through 0.02\% (at default) to 0.69\% (at the 75\% sweep point).

\item \textbf{Jump intensity $\lambda_J$} --- exogenous market condition. Higher jump frequency (12 vs.\ 4 per year) increases tail risk by amplifying crash-like events.

\item \textbf{Oracle decay $\alpha$} --- secondary effect. At the 33\% LTV floor, all decay rates produce zero bad debt. At the shipped 66.67\% default and above, faster decay reduces phantom-healthy windows, but the effect is dominated by the LTV choice.

\item \textbf{Refresh interval} --- operational parameter. Faster refreshes (15\,min vs.\ 2\,h) reduce the initial blind period but have diminishing returns once below the half-life.

\item \textbf{Lock fraction $\phi$} --- emergent, not directly controlled. Locks reduce cascade-induced bad debt but have minimal effect on oracle-staleness-driven bad debt.

\item \textbf{Liquidation-recovery haircut $\kappa$} --- exogenous. Increases bad debt during partial-liquidation steps in declining markets but does not affect oracle staleness. ($\kappa$ here is the dimensionless MC haircut of \S\ref{sec:bd-partial-liq}, not the linear-depth cascade coefficient $k$ of \S\ref{apx:cascade-sim}.)
\end{enumerate}

\subsubsection{Interaction Effects}
\label{subsec:bd-interactions}

Two notable interaction effects emerge from pairwise analysis:

\textbf{$\alpha \times$ LTV.} Higher LTV \emph{amplifies} the staleness problem non-linearly. At the 33\% LTV floor, changing $\alpha$ from 0.891 to 0.972 has zero observable effect; at the 75\% upper sweep point, the same change measurably increases bad debt. The over-collateralization buffer acts as a threshold: below a critical LTV, oracle delay is completely absorbed; the shipped 66.67\% default sits just above this threshold.

\textbf{$\alpha \times$ Refresh interval.} Diminishing returns to faster refreshes when $\alpha$ is already responsive. At $\alpha = 0.891$ (6\,h HL), reducing refresh interval from 2\,h to 15\,min has minimal effect because the EMA already converges rapidly.

% ==============================================================================
% 7. Partial Liquidation Model
% ==============================================================================
\subsection{Partial Liquidation}
\label{sec:bd-partial-liq}

The protocol uses a $2^{-e}$ partial liquidation model. \texttt{Pool.sol} computes \texttt{borrowed\_total >> partial\_exp} and \texttt{supplied\_total >> partial\_exp}, transferring fraction $2^{-e}$ of the position per liquidation call. Thus $e = 0$ corresponds to full liquidation (100\%), $e = 1$ clears 50\% per step (the default), $e = 2$ clears 25\%, and increasing $e$ liquidates a \emph{smaller} share each step, leaving more residual debt to be wound down across subsequent calls.

During continued price declines, partial liquidation can produce \emph{more} bad debt than instant full liquidation, because remaining position fractions face increasingly adverse prices. A dimensionless liquidation-recovery haircut $\kappa$ further increases bad debt by modeling the proceeds shortfall of forced liquidation sales:
\begin{equation}
p_{\text{effective}} = p_{\text{true}} \cdot (1 - \kappa)
\end{equation}
This $\kappa$ is distinct from the cascade-simulation depth coefficient $k$ of \S\ref{apx:cascade-sim}: $\kappa$ is a dimensionless fractional haircut applied to per-step recovery, whereas $k$ has units of $1/\text{supply}$ (linear-depth slippage of the cascade model). The two are not numerically comparable.

Monte Carlo tests confirm that a haircut $\kappa = 0.10$ (10\%) measurably increases bad debt relative to $\kappa = 0$ in declining markets. However, the effect is secondary to the LTV choice and primarily affects already-distressed positions.

% ==============================================================================
% 8. Cross-Validation
% ==============================================================================
\subsection{Cross-Validation}
\label{sec:bd-validation}

\subsubsection{Oracle Model vs.\ Solidity}
\label{subsec:bd-oracle-validation}

The Python oracle model is cross-validated against three Foundry scenario tests from \texttt{test/Oracle/OracleScenario.t.sol}:

\begin{itemize}[leftmargin=*]
\item \textbf{S01 (Benign Motion)}: 24 refreshes with $\pm$5\% oscillations. Mid drift $< 10\%$. \checkmark
\item \textbf{S02 (Sudden Drift)}: 2$\times$ price jump. After 1st refresh: $\delta < 3\%$; after 2nd: $\delta < 15\%$; after 14th: $15\% < \delta < 70\%$. \checkmark
\item \textbf{S04 (Sudden Reversal)}: 2$\times$ spike then revert. Residual $< 5\%$ (Solidity cadence) or $< 7\%$ (Python all-tick cadence). \checkmark
\end{itemize}

\noindent The discrepancy in S04 arises from Solidity's \texttt{delayed()} modifier using strict~$>$ comparison, causing only every other hourly tick to produce a Refresh event (alternating Pending/Refresh). Both models produce correct results for their respective refresh cadences.

\subsubsection{Phantom Window Verification}
\label{subsec:bd-phantom-validation}

Analytical phantom window computations (Section~\ref{subsec:bd-phantom}) are verified against step-by-step OracleModel simulations for 4 representative (crash, health) pairs. All match within $\pm 1$ refresh.

\subsubsection{Bound Conservatism}
\label{subsec:bd-bound-validation}

The analytical bound (Theorem~\ref{thm:bd-bound}) is verified to exceed all MC simulation results across all tested parameter combinations. The \texttt{check-bound} CLI command automates this check for each (LTV, $\alpha$, crash) triple in the MC results.

% ==============================================================================
% 9. Recommendations
% ==============================================================================
\subsection{Deployment Recommendations}
\label{sec:bd-recommendations}

Based on the analysis, we recommend the following parameter configurations for deployment:

\begin{enumerate}[leftmargin=*]
\item \textbf{Shipped default} ($\alpha = 0.944$, LTV = 66.67\%): Bounded bad debt of $\mathbb{E}[\text{BD}] = 0.00\%$ and CVaR(99\%) = 0.02\% under empirical jump-diffusion paths; analytical bound 1.90\% for 50\% crashes. This is the production setting, balancing capital efficiency against bad-debt containment.

\item \textbf{Conservative mode} ($\alpha = 0.944$, LTV = 33\%): Zero bounded bad debt for crashes up to 50\%---the governance-reachable floor for protocol-level conservatism, two lethargic cycles below the shipped default.

\item \textbf{Higher capital efficiency} ($\alpha = 0.891$, LTV = 50\%): Zero bounded bad debt for 50\% crashes. Requires reducing half-life to 6~hours, which weakens manipulation resistance (20~hours sustained manipulation for 90\% deviation).

\item \textbf{Maximum responsiveness} ($\alpha = 0.891$, LTV = 67\%): Bounded bad debt of 3.51\% for 50\% crash. Suitable only for highly liquid pairs where manipulation risk is low and faster oracle response is prioritized.
\end{enumerate}

\textbf{Configurations to avoid:}
\begin{itemize}[leftmargin=*]
\item LTV $\geq 75\%$ at any $\alpha$: Bounded bad debt exceeds 9\% for 50\% crashes.
\item $\alpha = 0.972$ (24\,h HL) with LTV $> 50\%$: Excessive staleness compounds with thin over-collateralization margin.
\end{itemize}

% ==============================================================================
% 10. Conclusion
% ==============================================================================
\subsection{Conclusion}
\label{sec:bd-conclusion}

XPower Banq's log-space TWAP oracle creates a quantifiable trade-off between manipulation resistance and bad-debt risk. At the shipped default of 66.67\% LTV, the 50\% over-collateralization buffer keeps the analytical and Monte Carlo bad-debt bounds tight: $\mathbb{E}[\text{BD}] = 0.00\%$, CVaR(99\%) = 0.02\%, and 1.3\% of jump-diffusion paths producing any bad debt. At the conservative-mode floor of 33\% LTV, the 200\% over-collateralization fully absorbs oracle staleness: Monte Carlo simulation produces zero bad debt across all 1{,}000 simulated paths, and the analytical bound confirms zero bad debt for crashes up to 50\%. Bad debt grows non-linearly above 67\%, reaching $\mathbb{E}[\text{BD}] = 0.01\%$ at the 75\% upper sweep point.

The mechanism underlying this robustness is the gap between the oracle's liquidation trigger ($p_{\text{crit}} = 1/H_0 \approx 0.67$, a 33\% decline) and the bad-debt threshold ($p = \text{LTV}/H_0$). At the conservative 33\% LTV floor, the threshold sits at $\approx 0.22$ (78\% decline)---a 45-percentage-point gap. At the shipped 66.67\% default, the threshold rises to $\approx 0.44$ (56\% decline), narrowing the gap to 23~percentage points but still preserving multi-hour blind-time tolerance. Even with phantom-healthy windows of 17--30~hours for 40--50\% crashes (Table~\ref{tab:phantom}), positions are liquidated long before the true price can traverse this gap. Only when the LTV approaches $1/H_0$---collapsing the gap---does oracle staleness translate into realised bad debt.

The primary risk lever is the LTV parameter, not the oracle decay rate. Sensitivity analysis (\S\ref{sec:bd-sensitivity}) ranks the six model parameters by impact: LTV dominates, followed by exogenous jump intensity $\lambda_J$, with oracle decay $\alpha$ a distant third. At the 33\% LTV floor, all three tested decay rates ($\alpha \in \{0.891, 0.944, 0.972\}$) produce identical zero bad debt (Table~\ref{tab:mc-alpha}), confirming that the over-collateralization buffer acts as a binary absorber below a critical LTV. Above that threshold (which the shipped 66.67\% default crosses by a small margin), the interaction between $\alpha$ and LTV becomes non-linear: increasing LTV toward 75\% amplifies the staleness penalty from $0\times$ to $1.67\times$ the instant-liquidation baseline.

Governance should therefore prioritize LTV conservatism over oracle responsiveness. The analytical bound $\text{BD}_{\max}(\delta, \alpha, \text{LTV}) \leq \text{BD}_{\text{inst}} \cdot (1 + W_{\max} \cdot \sigma_h / \delta)$ (Theorem~\ref{thm:bd-bound}) provides a closed-form, governance-usable formula for evaluating parameter changes against bad-debt tolerance. Combined with the safe-operating-region criterion ($\text{BD}_{\max} < 5\%$ of TVL for $\delta \leq 50\%$), this bound delineates five safe $(\alpha, \text{LTV})$ configurations (Table~\ref{tab:safe}), all with LTV $\leq 67\%$.

\textbf{Limitations.} The analysis makes several simplifying assumptions. First, positions are modelled with a static health distribution $H_0 \sim \mathcal{N}(1.5, 0.3)$ rather than endogenous dynamics from supply/borrow inflows. Second, oracle refreshes are assumed to occur at exact hourly intervals, whereas the \texttt{delayed()} modifier's strict-$>$ comparison can cause alternating Pending/Refresh patterns that effectively double the refresh interval (\S\ref{subsec:bd-oracle-validation}). Third, liquidation cascades and their feedback effects on market prices are modelled only via the liquidation-recovery haircut $\kappa$ (\S\ref{sec:bd-partial-liq}), not through an order-book simulation. Fourth, correlated multi-asset crashes are not considered.




% ==============================================================================
% Standalone-only References
% ==============================================================================
\ifbundle\else
\begin{thebibliography}{99}

\bibitem{wp}
\textsc{Karun The Ritch}.
\newblock XPower Banq: A Permissionless DeFi Lending Protocol with Lethargic Governance and Beta-Distributed Position Caps.
\newblock \emph{Preprint}, 2026.
\newblock \url{https://www.xpowerbanq.com}

\bibitem{mtp}
\textsc{Karun The Ritch}.
\newblock XPower Banq: Mathematical Theory \& Proofs.
\newblock \emph{Preprint}, 2026.
\newblock \url{https://www.xpowerbanq.com}

\bibitem{mackinga2022}
Mackinga, T., Nadahalli, T., and Wattenhofer, R.
\newblock TWAP Oracle Attacks: Easier Done than Said?
\newblock \emph{IEEE International Conference on Blockchain and Cryptocurrency}, 2022.

\bibitem{merton1976}
Merton, R.C.
\newblock Option Pricing When Underlying Stock Returns Are Discontinuous.
\newblock \emph{Journal of Financial Economics}, 3(1--2):125--144, 1976.

\bibitem{ref}
\textsc{Karun The Ritch}.
\newblock XPower Banq: References \& Glossary.
\newblock \emph{Preprint}, 2026.
\newblock \url{https://www.xpowerbanq.com}

\end{thebibliography}
\fi

% ==============================================================================
% Glossary
% ==============================================================================
\ifbundle\else
\section*{Glossary}
\label{sim:glossary}
\addcontentsline{toc}{section}{Glossary}

\noindent A comprehensive glossary with over 120 entries appears in \cite{ref}. The following terms are central to the bad-debt risk analysis in this paper; entries marked with a dagger~($\dagger$) denote symbols defined here.

\begin{description}[style=nextline, font=\normalfont\bfseries, leftmargin=1.5em]

\item[$A(n)$\,$\dagger$] Absorption fraction---the share of the log-space deviation absorbed after $n$ oracle refreshes: $A(n) = 1 - \alpha^{n-1}$ for $n \geq 2$, with $A(0) = A(1) = 0$ due to the one-refresh delay.

\item[Bad Debt] Uncollectable debt arising when a position's collateral value falls below its debt value; protocol absorbs the loss. Quantified in this paper.

\item[$\text{BD}_{\text{inst}}$\,$\dagger$] Bad debt under instant (zero-delay) liquidation. The baseline shortfall when a position's true collateral value $H_0(1{-}\delta)$ falls below LTV.

\item[$\text{BD}_{\max}$\,$\dagger$] Conservative upper bound on bad debt (Theorem~\ref{thm:bd-bound}), incorporating the oracle delay penalty: $\text{BD}_{\max} \leq \text{BD}_{\text{inst}} \cdot (1 + W_{\max}\,\sigma_h / \delta)$.

\item[$\delta$\,$\dagger$] Crash fraction. An instantaneous price drop from $p_0$ to $p_0(1{-}\delta)$; e.g.\ $\delta = 0.50$ is a 50\% crash.

\item[$\text{CVaR}(99\%)$\,$\dagger$] Conditional Value-at-Risk at the 99th percentile---the expected bad debt in the worst 1\% of simulated paths.

\item[$H_0$\,$\dagger$] Initial health factor of a position: $H_0 = (\text{supply} \times w_s) / (\text{borrow} \times w_b)$. Liquidation triggers when the oracle-observed health $H_{\text{oracle}} < 1$.

\item[Liquidation-Recovery Haircut ($\kappa$)\,$\dagger$] Fraction of collateral value lost to slippage and gas during liquidation; applied to the partial-liquidation recovery model.

\item[LTV] Loan-to-Value ratio---maximum borrowing power as fraction of collateral value; default $w_s/w_b = 170/255 \approx 66.67\%$.

\item[Merton Jump-Diffusion] Asset-price model with continuous Brownian motion plus a Poisson-driven jump component; basis for the Monte Carlo bad-debt simulation.

\item[Monte Carlo Simulation] Randomized numerical method using many simulated paths to estimate statistical distributions; used for bad-debt risk quantification and TWAP analysis (this paper).

\item[$p_{\text{crit}}$\,$\dagger$] Critical oracle price ratio that triggers liquidation: $p_{\text{crit}} = 1 / H_0$. The oracle fires when $\hat{p}(n)/p_0 < p_{\text{crit}}$.

\item[Phantom-healthy\,$\dagger$] A position state where the oracle reports $H_{\text{oracle}} \geq 1$ (solvent) but the true health $H_{\text{true}} < 1$ (underwater). Arises from oracle staleness during crashes (Definition~\ref{def:phantom}).

\item[$\sigma$, $\sigma_h$\,$\dagger$] Annualized and hourly volatility of the collateral asset. $\sigma_h = \sigma / \sqrt{8760}$. ETH calibration: $\sigma = 90\%$, $\sigma_h \approx 0.96\%$/hour.

\item[TVL\,$\dagger$] Total Value Locked---the aggregate collateral value in the lending pool. All bad-debt metrics are reported as a percentage of TVL.

\item[TWAP] Time-Weighted Average Price---price smoothed over time via log-space EMA (geometric mean temporal averaging) to resist manipulation. Simulated in this paper.

\item[$W$\,$\dagger$] Phantom-healthy window---the number of oracle refreshes during which a position remains phantom-healthy after a crash (Definition~\ref{def:phantom}).

\item[$\text{VaR}(99\%)$\,$\dagger$] Value-at-Risk at the 99th percentile---the bad debt level exceeded in only 1\% of simulated paths.

\end{description}

\clearpage
\ifodd\value{page}\else\thispagestyle{empty}\null\clearpage\fi
\fi

\end{document}
